% Expanded verification companion. Original physics sections are retained.
% Every table below is a transcription of the saved 113-field release.
\clearpage
\section{Weight-two interactions: a complete printed reconstruction atlas}
\label{expphysics:interaction-atlas}

The short physics chapter identifies the fields and quotes selected
couplings. This section supplies the information needed to inspect the
entire weight-two primary tensor without opening its 18,365-row sparse
file. It gives the actual component index dictionary, every highest-vector
norm, every nonzero ordered map between current multiplets, a
phase-sensitive representative of each map, and an explicit elementary
formula reconstructing the other components. The compression uses
current-algebra equivariance, not a numerical approximation or a choice
to omit small coefficients.

The source distinction matters. The files
\artifact{c8_physics_20260927/complete_primary_couplings.json} and its
neighboring TSV concern the earlier 96-field slice. All interaction
tables in this section use the files of the same names in
\artifact{c8_physics_20260927/full_weight_two/}. That later release
contains all 113 Virasoro primaries in $U_2$. Its exact sparse tensor has
$113^3=1,442,897$ ordered entries, of which 18,365 are nonzero and
1,424,532 vanish. The additional 17 fields are retained throughout this
section, including in the normalization and the $u$-pair sum rule.

\subsection{Actual component basis, including the lowering gauge}
\label{expphysics:lowering-gauge}

Write $V_m$ for the current multiplet numbered $m$, with dimensions
$(n_m,r_m)$ under the left and right $\mathfrak{sl}_2$ factors. The
field denoted by $(m,p,q)$ is the actual saved vector
\begin{equation}
 v_{m,p,q}=(F_L)_0^p(F_R)_0^qv_{m,0,0},
 \qquad 0\leq p<n_m,\quad0\leq q<r_m.
 \label{expphysics:component-basis}
\end{equation}
The integer indices $0,\ldots,112$ used in the tensor are printed in
Table~\ref{expphysics:basis-dictionary}. There is no factorial division
in this convention. The Cartan charges are
$(n_m-1-2p,r_m-1-2q)$, and the algebraic raising rules are
\begin{align}
 (E_L)_0v_{m,p,q}&=p(n_m-p)v_{m,p-1,q},\\
 (E_R)_0v_{m,p,q}&=q(r_m-q)v_{m,p,q-1}.
 \label{expphysics:raising}
\end{align}
These rules fix an algebraic $\mathfrak{sl}_2$ basis. They do not by
themselves say that the displayed $E$ and $F$ are Hermitian adjoints.

In fact the left lowering normalization is substantial and must not
be lost in a comparison of component tables. The saved positive norms
give
\begin{equation}
 \boxed{\displaystyle
 N_{m,p,q}=N_{m,0,0}\,8640^{2p}
 \frac{p!\,(n_m-1)!}{(n_m-1-p)!}
 \frac{q!\,(r_m-1)!}{(r_m-1-q)!}.}
 \label{expphysics:all-norms}
\end{equation}
Thus on these multiplets the adjoints of the zero-mode lowering
operators are $(F_L)_0^*=8640^2(E_L)_0$ and
$(F_R)_0^*=(E_R)_0$. Equivalently one can pass to compactly normalized
generators $E'_L=8640E_L$, $F'_L=F_L/8640$. The tabulated vectors
continue to use the original saved generators, as in
\eqref{expphysics:component-basis}. For example, the highest defining
field has $N_{7,0,0}=12$, whereas $N_{7,1,0}=895795200$.
Replacing the latter by 12 would change normalized coefficients
incorrectly. The right lowering alone has the familiar
factor $q!(r_m-1)!/(r_m-1-q)!$.

The factor can also be derived directly from the torus gauge,
without consulting all component norms. The left root is
$\alpha_L=(-1,-2,-3,-2,-1,-2,0,0)$ and the gauge entries are
$(5,8,9,8,5,1,2,2)$. Hence
\[
 \prod_i\widetilde c_i^{(\alpha_L)_i}
 =\frac1{5^2\,8^4\,9^3}=\frac1{8640^2}.
\]
With $E_L=e^{\alpha_L}$ and $F_L=-e^{-\alpha_L}$, transported
PCT gives $\theta_g(E_L)=-F_L/8640^2$.
The adjoint identity for a weight-one zero mode is
$a_0^*=-(\theta_ga)_0$, yielding the adjoints just stated.
The right-current coefficients similarly pair each root with its
inverse gauge factor, giving the unit compact normalization on
the right. Iterating these adjoints and the raising rules proves
\eqref{expphysics:all-norms}.

Equation~\eqref{expphysics:all-norms} is also directly visible in the
saved data: the transcription checks its ratio against every one of
the 113 recorded norms. This is a check of an encoding convention,
not a new lattice-mode calculation. The highest vectors and their
norms are as follows.

\begingroup\small
\setlength{\tabcolsep}{5pt}
\begin{longtable}{rrrrr}
\caption{All thirteen multiplet highest-vector norms in the saved lattice gauge. The last column counts stored nonzero Fock-basis terms; it is not a representation dimension.}\label{expphysics:highest-norms}\\
\toprule
$m$ & first index & $(n_m,r_m)$ & $N_{m,0,0}$ & Fock terms\\
\midrule
\endfirsthead
\multicolumn{5}{l}{\small\itshape Continued from preceding page.}\\
\toprule
$m$ & first index & $(n_m,r_m)$ & $N_{m,0,0}$ & Fock terms\\
\midrule
\endhead
\midrule
\multicolumn{5}{r}{\small\itshape Continued on next page.}\\
\endfoot
\bottomrule
\endlastfoot
0 & 0 & $( 1,5)$ & $\frac{20520}{91}$ & 61\\
1 & 5 & $( 1,9)$ & $3432$ & 32\\
2 & 14 & $( 1,9)$ & $3900$ & 32\\
3 & 23 & $( 1,13)$ & $\frac{1123200}{7}$ & 9\\
4 & 36 & $( 2,2)$ & $\frac{65}{135936}$ & 42\\
5 & 40 & $( 2,6)$ & $\frac{65}{2268}$ & 32\\
6 & 52 & $( 2,8)$ & $\frac{143}{752}$ & 22\\
7 & 68 & $( 2,14)$ & $12$ & 3\\
8 & 96 & $( 1,5)$ & $2964$ & 26\\
9 & 101 & $( 3,3)$ & $\frac{13}{24883200}$ & 7\\
10 & 110 & $( 1,1)$ & $\frac{9243}{1681}$ & 66\\
11 & 111 & $( 1,1)$ & $\frac{33748800}{1615441}$ & 76\\
12 & 112 & $( 1,1)$ & $\frac{49840}{961}$ & 74\\
\end{longtable}
\endgroup


The first eight highest vectors are the original charged multiplets,
of types $(1,5)$, $(1,9)$, $(1,9)$, $(1,13)$, $(2,2)$,
$(2,6)$, $(2,8)$ and $(2,14)$.
The two copies of $(1,9)$ have different coset weights; their labels
must remain distinct. Multiplets 8 and 9 have highest vectors
$(E_R)_{-1}E_R$ and $(E_L)_{-1}E_R$ and types $(1,5),(3,3)$.
The last three are the following total-Virasoro primaries:
\begin{align}
 v_{10}&=c_rT_L-c_lT_R,\\
 v_{11}&=c_e(T_L+T_R)-(c_l+c_r)T_e,\\
 v_{12}&=c_f(T_L+T_R+T_e)-(c_l+c_r+c_e)T_f,
 \label{expphysics:stress-differences}
\end{align}
where
\[
 (c_l,c_r,c_e,c_f)=
 \left(1,\frac{117}{41},\frac{2400}{1271},\frac{70}{31}\right).
\]
The stress labels here are internal commuting chiral factors, not
the left and right spacetime central charges. Their sum is eight.

For a linear combination $A=\sum_i a_iT_i$, the fourth-order pole of
the total stress with $A$ is $\tfrac12\sum_i a_ic_i$. Hence $A$ is
primary of weight two exactly when $\sum_i a_ic_i=0$. Since the four
stress tensors are orthogonal and have norms $c_i/2$, the three
combinations above are mutually orthogonal and orthogonal to the
total stress. This verifies both the primary condition and the
normalization mechanism without referring to a mode engine.

There are $96+5+9+3=113$ displayed fields. The vacuum character gives
$\dim U_2=120$ and $\dim U_1=6$. The six current derivatives and the
total stress account for the remaining seven states. The positive
orthogonal primary vectors therefore exhaust the primary space:
there is no omitted primary multiplet at this weight. The first
column of the next table makes all subsequent tensor indices
unambiguous.

\begingroup\small
\setlength{\tabcolsep}{5pt}
\begin{longtable}{rrrr}
\caption{Complete index dictionary for all 113 weight-two Virasoro primaries. Multiplying the last column by Table~\ref{expphysics:highest-norms} gives the actual squared norm.}\label{expphysics:basis-dictionary}\\
\toprule
index $i$ & $(m,p,q)$ & $(H_L,H_R)$ & $N_i/N_{m,0,0}$\\
\midrule
\endfirsthead
\multicolumn{4}{l}{\small\itshape Continued from preceding page.}\\
\toprule
index $i$ & $(m,p,q)$ & $(H_L,H_R)$ & $N_i/N_{m,0,0}$\\
\midrule
\endhead
\midrule
\multicolumn{4}{r}{\small\itshape Continued on next page.}\\
\endfoot
\bottomrule
\endlastfoot
0 & $(0,0,0)$ & $(0,4)$ & 1\\
1 & $(0,0,1)$ & $(0,2)$ & 4\\
2 & $(0,0,2)$ & $(0,0)$ & 24\\
3 & $(0,0,3)$ & $(0,-2)$ & 144\\
4 & $(0,0,4)$ & $(0,-4)$ & 576\\
5 & $(1,0,0)$ & $(0,8)$ & 1\\
6 & $(1,0,1)$ & $(0,6)$ & 8\\
7 & $(1,0,2)$ & $(0,4)$ & 112\\
8 & $(1,0,3)$ & $(0,2)$ & 2016\\
9 & $(1,0,4)$ & $(0,0)$ & 40320\\
10 & $(1,0,5)$ & $(0,-2)$ & 806400\\
11 & $(1,0,6)$ & $(0,-4)$ & 14515200\\
12 & $(1,0,7)$ & $(0,-6)$ & 203212800\\
13 & $(1,0,8)$ & $(0,-8)$ & 1625702400\\
14 & $(2,0,0)$ & $(0,8)$ & 1\\
15 & $(2,0,1)$ & $(0,6)$ & 8\\
16 & $(2,0,2)$ & $(0,4)$ & 112\\
17 & $(2,0,3)$ & $(0,2)$ & 2016\\
18 & $(2,0,4)$ & $(0,0)$ & 40320\\
19 & $(2,0,5)$ & $(0,-2)$ & 806400\\
20 & $(2,0,6)$ & $(0,-4)$ & 14515200\\
21 & $(2,0,7)$ & $(0,-6)$ & 203212800\\
22 & $(2,0,8)$ & $(0,-8)$ & 1625702400\\
23 & $(3,0,0)$ & $(0,12)$ & 1\\
24 & $(3,0,1)$ & $(0,10)$ & 12\\
25 & $(3,0,2)$ & $(0,8)$ & 264\\
26 & $(3,0,3)$ & $(0,6)$ & 7920\\
27 & $(3,0,4)$ & $(0,4)$ & 285120\\
28 & $(3,0,5)$ & $(0,2)$ & 11404800\\
29 & $(3,0,6)$ & $(0,0)$ & 479001600\\
30 & $(3,0,7)$ & $(0,-2)$ & 20118067200\\
31 & $(3,0,8)$ & $(0,-4)$ & 804722688000\\
32 & $(3,0,9)$ & $(0,-6)$ & 28970016768000\\
33 & $(3,0,10)$ & $(0,-8)$ & 869100503040000\\
34 & $(3,0,11)$ & $(0,-10)$ & 19120211066880000\\
35 & $(3,0,12)$ & $(0,-12)$ & 229442532802560000\\
36 & $(4,0,0)$ & $(1,1)$ & 1\\
37 & $(4,0,1)$ & $(1,-1)$ & 1\\
38 & $(4,1,0)$ & $(-1,1)$ & 74649600\\
39 & $(4,1,1)$ & $(-1,-1)$ & 74649600\\
40 & $(5,0,0)$ & $(1,5)$ & 1\\
41 & $(5,0,1)$ & $(1,3)$ & 5\\
42 & $(5,0,2)$ & $(1,1)$ & 40\\
43 & $(5,0,3)$ & $(1,-1)$ & 360\\
44 & $(5,0,4)$ & $(1,-3)$ & 2880\\
45 & $(5,0,5)$ & $(1,-5)$ & 14400\\
46 & $(5,1,0)$ & $(-1,5)$ & 74649600\\
47 & $(5,1,1)$ & $(-1,3)$ & 373248000\\
48 & $(5,1,2)$ & $(-1,1)$ & 2985984000\\
49 & $(5,1,3)$ & $(-1,-1)$ & 26873856000\\
50 & $(5,1,4)$ & $(-1,-3)$ & 214990848000\\
51 & $(5,1,5)$ & $(-1,-5)$ & 1074954240000\\
52 & $(6,0,0)$ & $(1,7)$ & 1\\
53 & $(6,0,1)$ & $(1,5)$ & 7\\
54 & $(6,0,2)$ & $(1,3)$ & 84\\
55 & $(6,0,3)$ & $(1,1)$ & 1260\\
56 & $(6,0,4)$ & $(1,-1)$ & 20160\\
57 & $(6,0,5)$ & $(1,-3)$ & 302400\\
58 & $(6,0,6)$ & $(1,-5)$ & 3628800\\
59 & $(6,0,7)$ & $(1,-7)$ & 25401600\\
60 & $(6,1,0)$ & $(-1,7)$ & 74649600\\
61 & $(6,1,1)$ & $(-1,5)$ & 522547200\\
62 & $(6,1,2)$ & $(-1,3)$ & 6270566400\\
63 & $(6,1,3)$ & $(-1,1)$ & 94058496000\\
64 & $(6,1,4)$ & $(-1,-1)$ & 1504935936000\\
65 & $(6,1,5)$ & $(-1,-3)$ & 22574039040000\\
66 & $(6,1,6)$ & $(-1,-5)$ & 270888468480000\\
67 & $(6,1,7)$ & $(-1,-7)$ & 1896219279360000\\
68 & $(7,0,0)$ & $(1,13)$ & 1\\
69 & $(7,0,1)$ & $(1,11)$ & 13\\
70 & $(7,0,2)$ & $(1,9)$ & 312\\
71 & $(7,0,3)$ & $(1,7)$ & 10296\\
72 & $(7,0,4)$ & $(1,5)$ & 411840\\
73 & $(7,0,5)$ & $(1,3)$ & 18532800\\
74 & $(7,0,6)$ & $(1,1)$ & 889574400\\
75 & $(7,0,7)$ & $(1,-1)$ & 43589145600\\
76 & $(7,0,8)$ & $(1,-3)$ & 2092278988800\\
77 & $(7,0,9)$ & $(1,-5)$ & 94152554496000\\
78 & $(7,0,10)$ & $(1,-7)$ & 3766102179840000\\
79 & $(7,0,11)$ & $(1,-9)$ & 124281371934720000\\
80 & $(7,0,12)$ & $(1,-11)$ & 2982752926433280000\\
81 & $(7,0,13)$ & $(1,-13)$ & 38775788043632640000\\
82 & $(7,1,0)$ & $(-1,13)$ & 74649600\\
83 & $(7,1,1)$ & $(-1,11)$ & 970444800\\
84 & $(7,1,2)$ & $(-1,9)$ & 23290675200\\
85 & $(7,1,3)$ & $(-1,7)$ & 768592281600\\
86 & $(7,1,4)$ & $(-1,5)$ & 30743691264000\\
87 & $(7,1,5)$ & $(-1,3)$ & 1383466106880000\\
88 & $(7,1,6)$ & $(-1,1)$ & 66406373130240000\\
89 & $(7,1,7)$ & $(-1,-1)$ & 3253912283381760000\\
90 & $(7,1,8)$ & $(-1,-3)$ & 156187789602324480000\\
91 & $(7,1,9)$ & $(-1,-5)$ & 7028450532104601600000\\
92 & $(7,1,10)$ & $(-1,-7)$ & 281138021284184064000000\\
93 & $(7,1,11)$ & $(-1,-9)$ & 9277554702378074112000000\\
94 & $(7,1,12)$ & $(-1,-11)$ & 222661312857073778688000000\\
95 & $(7,1,13)$ & $(-1,-13)$ & 2894597067141959122944000000\\
96 & $(8,0,0)$ & $(0,4)$ & 1\\
97 & $(8,0,1)$ & $(0,2)$ & 4\\
98 & $(8,0,2)$ & $(0,0)$ & 24\\
99 & $(8,0,3)$ & $(0,-2)$ & 144\\
100 & $(8,0,4)$ & $(0,-4)$ & 576\\
101 & $(9,0,0)$ & $(2,2)$ & 1\\
102 & $(9,0,1)$ & $(2,0)$ & 2\\
103 & $(9,0,2)$ & $(2,-2)$ & 4\\
104 & $(9,1,0)$ & $(0,2)$ & 149299200\\
105 & $(9,1,1)$ & $(0,0)$ & 298598400\\
106 & $(9,1,2)$ & $(0,-2)$ & 597196800\\
107 & $(9,2,0)$ & $(-2,2)$ & 22290251120640000\\
108 & $(9,2,1)$ & $(-2,0)$ & 44580502241280000\\
109 & $(9,2,2)$ & $(-2,-2)$ & 89161004482560000\\
110 & $(10,0,0)$ & $(0,0)$ & 1\\
111 & $(11,0,0)$ & $(0,0)$ & 1\\
112 & $(12,0,0)$ & $(0,0)$ & 1\\
\end{longtable}
\endgroup


\subsection{What a phase-sensitive representative means}
\label{expphysics:phase-convention}

Let $N_i=\langle v_i,v_i\rangle_g>0$ and
$\widehat v_i=v_i/\sqrt{N_i}$, using the positive square root.
All pairings are antilinear in their first slot. The tensor is
\begin{equation}
 C_{ij}^{\ k}=
 \frac{\langle v_k,(v_i)_1v_j\rangle_g}
      {\sqrt{N_iN_jN_k}}.
 \label{expphysics:normalized-tensor}
\end{equation}
Thus its third entry refers to the adjoint field at infinity. In a
local OPE it multiplies $z^{-2}\widehat v_k$. Projection to the
primary space is understood: the identity, currents, their
derivatives, and the total stress have their own OPE terms.

To print complex phases compactly, use only in this section
\[
 \sigma=\sqrt{13}>0,\qquad\vartheta=i\sqrt3,
 \qquad\overline\sigma=\sigma,\quad
 \overline\vartheta=-\vartheta.
\]
The numerator in \eqref{expphysics:normalized-tensor} is a positive
rational multiple of a primitive integer combination
$\kappa=a+b\sigma+c\vartheta+d\sigma\vartheta$. ``Primitive''
means that the four integer coefficients have no common positive
factor. Its sign is retained; $\kappa$ and $-\kappa$ are different
phases. Each representative table row gives $\kappa$ and the saved
rational number $Q=|C_{ij}^{\ k}|^2$. The exact coefficient is
\begin{equation}
 \boxed{\displaystyle
 C_{ij}^{\ k}=\kappa
          \sqrt{\frac{Q}{\overline\kappa\kappa}}.}
 \label{expphysics:phase-reading}
\end{equation}
Every square root in this formula is positive real. This is a
complete exact complex value, including its phase; it is not merely
an absolute-value table. If desired its phase denominator is obtained
with rational arithmetic from
\[
 \overline\kappa\kappa=
 a^2+13b^2+3c^2+39d^2+(2ab+6cd)\sigma.
\]
For $\kappa=\pm1$ the coefficient is simply $\pm\sqrt Q$.
The use of $\sigma$ here avoids confusing $\sqrt{13}$ with the
dimension symbol $s=3d+2$ used in the modular-data discussion.

\subsection{An explicit reconstruction formula for every component}
\label{expphysics:transvectant}

For each fixed ordered triple $(m,n,\ell)$ the projection
\[
 P_\ell\bigl((\,\cdot\,)_1\,\cdot\,\bigr):V_m\otimes V_n\longrightarrow V_\ell
\]
is an $\mathfrak{sl}_2\oplus\mathfrak{sl}_2$ intertwiner. This
follows because current zero modes are derivations and commute with
the total Virasoro modes. Orthogonal projection onto the full
primary space preserves this property. Each factor in a finite
$\mathfrak{sl}_2$ tensor product occurs with multiplicity at most
one. Consequently a specified triple of the labeled current
multiplets has at most one intertwiner, even when two different
multiplets elsewhere in the list have the same dimensions.
One nonzero component fixes this map completely.

Here is a concrete formula that removes the need to choose a
Clebsch--Gordan phase convention. For one current factor let $a,b,c$
be the highest Cartan weights, i.e. the representation dimensions
minus one, and set $L=(a+b-c)/2$. A map is allowed only if
\[
 L\in\mathbb Z,\qquad0\leq L\leq\min(a,b).
\]
Use the polynomial realization with $F=Y\partial_X$ and basis
\[
 w_{a,p}=F^pX^a=a^{\underline p}X^{a-p}Y^p,
 \qquad a^{\underline p}=\frac{a!}{(a-p)!}.
\]
The $L$th transvectant is
\begin{equation}
 \mathcal T_L(f,g)=\sum_{t=0}^{L}(-1)^t\binom Lt
 (\partial_X^{L-t}\partial_Y^tf)
 (\partial_X^t\partial_Y^{L-t}g).
 \label{expphysics:transvectant-definition}
\end{equation}
It is equivariant and lands in degree $c=a+b-2L$. Put $r=p+q-L$.
The coefficient of $w_{c,r}$ in
$\mathcal T_L(w_{a,p},w_{b,q})$ is
\begin{equation}
 \begin{split}
 K_{a,b;c}(p,q;r)
 ={}&\frac{a^{\underline p}b^{\underline q}}{c^{\underline r}}
 \sum_{t=0}^{L}(-1)^t\binom Lt
 (a-p)^{\underline{L-t}}p^{\underline t}\,
 (b-q)^{\underline t}q^{\underline{L-t}}.
 \end{split}
 \label{expphysics:K-formula}
\end{equation}
Set this coefficient to zero if $r\ne p+q-L$ or $r$ is outside
$0,\ldots,c$. A falling factorial with more differentiations than
its nonnegative integer degree is zero. For degree zero,
$0^{\underline0}=1$. These conventions cover trivial representations,
tops and bottoms, without a limiting prescription.

For a full component triple $(i,j,k)$, multiply the $K$ factors
for the left and right current algebras and denote the result by
$K_{ijk}$. Use the lowering exponents from
Table~\ref{expphysics:basis-dictionary}. If $(i_*,j_*,k_*)$ is the
printed nonzero representative for the same ordered multiplet
triple, the exact normalized tensor is
\begin{equation}
 \boxed{\displaystyle
 C_{ij}^{\ k}=C_{i_*j_*}^{\ k_*}
 \frac{K_{ijk}}{K_{i_*j_*k_*}}
 \sqrt{\frac{N_kN_{i_*}N_{j_*}}
              {N_iN_jN_{k_*}}}.}
 \label{expphysics:full-reconstruction}
\end{equation}
The denominator $K_{i_*j_*k_*}$ is nonzero because that component
of the unique intertwiner is nonzero. All norm factors are explicit
in \eqref{expphysics:all-norms} and the printed highest-norm table.
The formula uses the raw algebraic lowering basis, so the factor
$8640^{2p}$ must be included in its norm conversion.

To see the normalization directly, write
$P_\ell((v_i)_1v_j)=B_{mn}^{\ \ell}K_{ijk}v_k$ in a
charge-compatible component. Pairing with $v_k$ gives a numerator
$B_{mn}^{\ \ell}K_{ijk}N_k$, hence
$C_{ij}^{\ k}=B_{mn}^{\ \ell}K_{ijk}\sqrt{N_k/(N_iN_j)}$.
Taking the ratio with the representative gives
\eqref{expphysics:full-reconstruction}; the unknown reduced scalar
$B_{mn}^{\ \ell}$ cancels. This proof also explains why the phase
of a representative is sufficient.

For example the first $V_0\otimes V_0\to V_1$ representative is
$(i,j,k)=(0,0,5)$, with
\[
 C_{00}^{\ 5}=\frac{\sigma+2\vartheta}{5}
                  \sqrt{\frac{3971}{2808}}.
\]
On the right factor its degrees are $(4,4,8)$, so $L=0$.
Moving to $(i,j,k)=(0,1,6)$ changes $K$ from 1 to $4/8=1/2$,
while $N_6/N_5=8$ and $N_1/N_0=4$. The normalized coefficient is
therefore divided by $\sqrt2$, and
$|C_{01}^{\ 6}|^2=3971/5616$, agreeing with the component export.
This example tests both the no-factorial convention and the
distinction between raw coefficients and unit-field coefficients.

The saved computation used a second, independent way to obtain
the components: the Ward recurrence
\[
 P((Fv_i)_1v_j)=F\,P((v_i)_1v_j)-P((v_i)_1Fv_j).
\]
Its directly computed highest-first slice contained 1137 actual
lattice mode products and 3512 charge-compatible projections,
with 1639 nonzero projections. After completion, the saved report
records 51,076 raising/lowering pair checks, 12,769 primary
skew-symmetry pair checks, and 33 additional direct lattice product
probes. The printed reconstruction formula is the representation
theory underlying those Ward identities; the document transcription
does not rerun the lattice calculation or extend its weight cutoff.

\subsection{Exhaustive support and all reduced coefficients}
\label{expphysics:atlas-coverage}

Among the $13^3=2197$ ordered multiplet triples, exactly 318 have
a nonzero projected map. The other 1879 maps vanish in the saved
complete tensor. This includes triples forbidden by the two
$\mathfrak{sl}_2$ triangle rules and any allowed triples whose reduced
coefficient is zero. A zero map is stronger than a vanishing selected
component. To make the distinction explicit, the support table lists
every ordered pair $(m,n)$, including pairs with no target.

\begingroup\small
\setlength{\tabcolsep}{5pt}
\begin{longtable}{rrlr}
\caption{Exhaustive support of the primary tensor by ordered pair of multiplets. A missing target means the entire equivariant map vanishes; the final column counts all nonzero ordered component triples for that pair.}\label{expphysics:coupling-support}\\
\toprule
first $m$ & second $n$ & nonzero target multiplets & nonzero components\\
\midrule
\endfirsthead
\multicolumn{4}{l}{\small\itshape Continued from preceding page.}\\
\toprule
first $m$ & second $n$ & nonzero target multiplets & nonzero components\\
\midrule
\endhead
\midrule
\multicolumn{4}{r}{\small\itshape Continued on next page.}\\
\endfoot
\bottomrule
\endlastfoot
0 & 0 & $0,1,2,8,10,11,12$ & 103\\
0 & 1 & $0,1,2,3$ & 148\\
0 & 2 & $0,1,2,3$ & 148\\
0 & 3 & $1,2,3$ & 149\\
0 & 4 & $5$ & 20\\
0 & 5 & $4,5,6$ & 120\\
0 & 6 & $5,6$ & 120\\
0 & 7 & $7$ & 124\\
0 & 8 & $0$ & 19\\
0 & 9 & $\varnothing$ & 0\\
0 & 10 & $0$ & 5\\
0 & 11 & $0$ & 5\\
0 & 12 & $0$ & 5\\
1 & 0 & $0,1,2,3$ & 148\\
1 & 1 & $0,1,2,3,8,10,11,12$ & 298\\
1 & 2 & $0,1,2,3$ & 232\\
1 & 3 & $0,1,2,3$ & 284\\
1 & 4 & $6$ & 32\\
1 & 5 & $5,6,7$ & 252\\
1 & 6 & $4,5,6,7$ & 352\\
1 & 7 & $5,6,7$ & 452\\
1 & 8 & $1$ & 39\\
1 & 9 & $\varnothing$ & 0\\
1 & 10 & $1$ & 9\\
1 & 11 & $1$ & 9\\
1 & 12 & $1$ & 9\\
2 & 0 & $0,1,2,3$ & 148\\
2 & 1 & $0,1,2,3$ & 232\\
2 & 2 & $0,1,2,3,8,10,11,12$ & 298\\
2 & 3 & $0,1,2,3$ & 284\\
2 & 4 & $6$ & 32\\
2 & 5 & $5,6,7$ & 252\\
2 & 6 & $4,5,6,7$ & 352\\
2 & 7 & $5,6,7$ & 452\\
2 & 8 & $2$ & 39\\
2 & 9 & $\varnothing$ & 0\\
2 & 10 & $2$ & 9\\
2 & 11 & $2$ & 9\\
2 & 12 & $2$ & 9\\
3 & 0 & $1,2,3$ & 149\\
3 & 1 & $0,1,2,3$ & 284\\
3 & 2 & $0,1,2,3$ & 284\\
3 & 3 & $0,1,2,3,8,10,11,12$ & 478\\
3 & 4 & $7$ & 52\\
3 & 5 & $6,7$ & 240\\
3 & 6 & $5,6,7$ & 392\\
3 & 7 & $4,5,6,7$ & 648\\
3 & 8 & $3$ & 59\\
3 & 9 & $\varnothing$ & 0\\
3 & 10 & $3$ & 13\\
3 & 11 & $3$ & 13\\
3 & 12 & $3$ & 13\\
4 & 0 & $5$ & 20\\
4 & 1 & $6$ & 32\\
4 & 2 & $6$ & 32\\
4 & 3 & $7$ & 52\\
4 & 4 & $9,10,11,12$ & 28\\
4 & 5 & $0$ & 20\\
4 & 6 & $1,2$ & 64\\
4 & 7 & $3$ & 52\\
4 & 8 & $\varnothing$ & 0\\
4 & 9 & $4$ & 16\\
4 & 10 & $4$ & 4\\
4 & 11 & $4$ & 4\\
4 & 12 & $4$ & 4\\
5 & 0 & $4,5,6$ & 120\\
5 & 1 & $5,6,7$ & 252\\
5 & 2 & $5,6,7$ & 252\\
5 & 3 & $6,7$ & 240\\
5 & 4 & $0$ & 20\\
5 & 5 & $0,1,2,8,9,10,11,12$ & 308\\
5 & 6 & $0,1,2,3$ & 320\\
5 & 7 & $1,2,3$ & 360\\
5 & 8 & $5$ & 44\\
5 & 9 & $5$ & 64\\
5 & 10 & $5$ & 12\\
5 & 11 & $5$ & 12\\
5 & 12 & $5$ & 12\\
6 & 0 & $5,6$ & 120\\
6 & 1 & $4,5,6,7$ & 352\\
6 & 2 & $4,5,6,7$ & 352\\
6 & 3 & $5,6,7$ & 392\\
6 & 4 & $1,2$ & 64\\
6 & 5 & $0,1,2,3$ & 320\\
6 & 6 & $0,1,2,3,8,9,10,11,12$ & 568\\
6 & 7 & $1,2,3$ & 464\\
6 & 8 & $6$ & 64\\
6 & 9 & $6$ & 88\\
6 & 10 & $6$ & 16\\
6 & 11 & $6$ & 16\\
6 & 12 & $6$ & 16\\
7 & 0 & $7$ & 124\\
7 & 1 & $5,6,7$ & 452\\
7 & 2 & $5,6,7$ & 452\\
7 & 3 & $4,5,6,7$ & 648\\
7 & 4 & $3$ & 52\\
7 & 5 & $1,2,3$ & 360\\
7 & 6 & $1,2,3$ & 464\\
7 & 7 & $0,1,2,3,8,9,10,11,12$ & 1168\\
7 & 8 & $7$ & 124\\
7 & 9 & $7$ & 160\\
7 & 10 & $7$ & 28\\
7 & 11 & $7$ & 28\\
7 & 12 & $7$ & 28\\
8 & 0 & $0$ & 19\\
8 & 1 & $1$ & 39\\
8 & 2 & $2$ & 39\\
8 & 3 & $3$ & 59\\
8 & 4 & $\varnothing$ & 0\\
8 & 5 & $5$ & 44\\
8 & 6 & $6$ & 64\\
8 & 7 & $7$ & 124\\
8 & 8 & $8,10,11,12$ & 34\\
8 & 9 & $9$ & 27\\
8 & 10 & $8$ & 5\\
8 & 11 & $8$ & 5\\
8 & 12 & $8$ & 5\\
9 & 0 & $\varnothing$ & 0\\
9 & 1 & $\varnothing$ & 0\\
9 & 2 & $\varnothing$ & 0\\
9 & 3 & $\varnothing$ & 0\\
9 & 4 & $4$ & 16\\
9 & 5 & $5$ & 64\\
9 & 6 & $6$ & 88\\
9 & 7 & $7$ & 160\\
9 & 8 & $9$ & 27\\
9 & 9 & $8,9,10,11,12$ & 90\\
9 & 10 & $9$ & 9\\
9 & 11 & $9$ & 9\\
9 & 12 & $9$ & 9\\
10 & 0 & $0$ & 5\\
10 & 1 & $1$ & 9\\
10 & 2 & $2$ & 9\\
10 & 3 & $3$ & 13\\
10 & 4 & $4$ & 4\\
10 & 5 & $5$ & 12\\
10 & 6 & $6$ & 16\\
10 & 7 & $7$ & 28\\
10 & 8 & $8$ & 5\\
10 & 9 & $9$ & 9\\
10 & 10 & $10,11,12$ & 3\\
10 & 11 & $10$ & 1\\
10 & 12 & $10$ & 1\\
11 & 0 & $0$ & 5\\
11 & 1 & $1$ & 9\\
11 & 2 & $2$ & 9\\
11 & 3 & $3$ & 13\\
11 & 4 & $4$ & 4\\
11 & 5 & $5$ & 12\\
11 & 6 & $6$ & 16\\
11 & 7 & $7$ & 28\\
11 & 8 & $8$ & 5\\
11 & 9 & $9$ & 9\\
11 & 10 & $10$ & 1\\
11 & 11 & $11,12$ & 2\\
11 & 12 & $11$ & 1\\
12 & 0 & $0$ & 5\\
12 & 1 & $1$ & 9\\
12 & 2 & $2$ & 9\\
12 & 3 & $3$ & 13\\
12 & 4 & $4$ & 4\\
12 & 5 & $5$ & 12\\
12 & 6 & $6$ & 16\\
12 & 7 & $7$ & 28\\
12 & 8 & $8$ & 5\\
12 & 9 & $9$ & 9\\
12 & 10 & $10$ & 1\\
12 & 11 & $11$ & 1\\
12 & 12 & $12$ & 1\\
\end{longtable}
\endgroup


The next table contains exactly one row for each of the 318 nonzero
maps. The choice of representative favors a short saved squared
coefficient and need not put the first component at the top of its
multiplet. Its first column identifies the ordered map; its second
column gives the actual component indices. These are all the rows
required by \eqref{expphysics:full-reconstruction}. A triple absent
from this table is a zero map, as specified by the support table.

\begingroup\small
\setlength{\tabcolsep}{5pt}
\begin{longtable}{rrlr}
\caption{Complete phase-sensitive representative atlas: one nonzero component for every nonzero ordered multiplet map. The coefficient is $C_{ij}^{\ k}=\kappa\sqrt{Q/(\overline\kappa\kappa)}$, with positive square roots. Use the index dictionary and the transvectant reconstruction in the text.}\label{expphysics:coupling-atlas}\\
\toprule
$(m,n,\ell)$ & $(i,j,k)$ & $\kappa$ & $Q=|C_{ij}^{\ k}|^2$\\
\midrule
\endfirsthead
\multicolumn{4}{l}{\small\itshape Continued from preceding page.}\\
\toprule
$(m,n,\ell)$ & $(i,j,k)$ & $\kappa$ & $Q=|C_{ij}^{\ k}|^2$\\
\midrule
\endhead
\midrule
\multicolumn{4}{r}{\small\itshape Continued on next page.}\\
\endfoot
\bottomrule
\endlastfoot
$(0,0,0)$ & $(0,2,0)$ & $1$ & $\frac{4356}{8645}$\\
$(0,0,1)$ & $(0,0,5)$ & $\sigma+2\vartheta$ & $\frac{3971}{2808}$\\
$(0,0,2)$ & $(0,0,14)$ & $-\sigma+3\vartheta$ & $\frac{361}{1560}$\\
$(0,0,8)$ & $(0,2,96)$ & $1$ & $\frac{2}{247}$\\
$(0,0,10)$ & $(0,4,110)$ & $-1$ & $\frac{4}{1027}$\\
$(0,0,11)$ & $(0,4,111)$ & $-1$ & $\frac{154587}{703100}$\\
$(0,0,12)$ & $(0,4,112)$ & $-1$ & $\frac{1681}{3115}$\\
$(0,1,0)$ & $(0,12,3)$ & $\sigma-2\vartheta$ & $\frac{3971}{5616}$\\
$(0,1,1)$ & $(0,7,5)$ & $1$ & $\frac{39605}{717288}$\\
$(0,1,2)$ & $(0,7,14)$ & $1+\sigma\vartheta$ & $\frac{57}{1144}$\\
$(0,1,3)$ & $(0,5,23)$ & $6-2\sigma+4\vartheta+\sigma\vartheta$ & $\frac{665}{858}$\\
$(0,2,0)$ & $(0,20,2)$ & $-\sigma-3\vartheta$ & $\frac{361}{7280}$\\
$(0,2,1)$ & $(0,16,5)$ & $1-\sigma\vartheta$ & $\frac{57}{1144}$\\
$(0,2,2)$ & $(0,16,14)$ & $-1$ & $\frac{39}{760}$\\
$(0,2,3)$ & $(0,15,24)$ & $9+2\sigma+6\vartheta-\sigma\vartheta$ & $\frac{19}{91}$\\
$(0,3,1)$ & $(0,35,13)$ & $6-2\sigma-4\vartheta-\sigma\vartheta$ & $\frac{665}{858}$\\
$(0,3,2)$ & $(0,34,21)$ & $9+2\sigma-6\vartheta+\sigma\vartheta$ & $\frac{19}{91}$\\
$(0,3,3)$ & $(2,24,24)$ & $1$ & $\frac{80}{1729}$\\
$(0,4,5)$ & $(0,36,40)$ & $-16-\sigma\vartheta$ & $\frac{19}{78}$\\
$(0,5,4)$ & $(0,45,37)$ & $-16+\sigma\vartheta$ & $\frac{19}{78}$\\
$(0,5,5)$ & $(1,42,41)$ & $1$ & $\frac{1815}{84721}$\\
$(0,5,6)$ & $(1,43,55)$ & $81-11\sigma\vartheta$ & $\frac{836}{4095}$\\
$(0,6,5)$ & $(1,56,42)$ & $81+11\sigma\vartheta$ & $\frac{836}{4095}$\\
$(0,6,6)$ & $(2,52,52)$ & $1$ & $\frac{208}{5985}$\\
$(0,7,7)$ & $(0,70,68)$ & $-1$ & $\frac{5}{133}$\\
$(0,8,0)$ & $(0,98,0)$ & $1$ & $\frac{2}{247}$\\
$(0,10,0)$ & $(0,110,0)$ & $-1$ & $\frac{4}{1027}$\\
$(0,11,0)$ & $(0,111,0)$ & $-1$ & $\frac{154587}{703100}$\\
$(0,12,0)$ & $(0,112,0)$ & $-1$ & $\frac{1681}{3115}$\\
$(1,0,0)$ & $(5,4,0)$ & $\sigma-2\vartheta$ & $\frac{3971}{2808}$\\
$(1,0,1)$ & $(5,4,7)$ & $1$ & $\frac{39605}{717288}$\\
$(1,0,2)$ & $(5,2,14)$ & $1+\sigma\vartheta$ & $\frac{133}{572}$\\
$(1,0,3)$ & $(5,0,23)$ & $6-2\sigma+4\vartheta+\sigma\vartheta$ & $\frac{665}{858}$\\
$(1,1,0)$ & $(5,11,0)$ & $1-4\sigma\vartheta$ & $\frac{39605}{717288}$\\
$(1,1,1)$ & $(7,9,7)$ & $-\sigma+2\vartheta$ & $\frac{5041}{960960}$\\
$(1,1,2)$ & $(7,9,16)$ & $-7\sigma-11\vartheta$ & $\frac{25}{17472}$\\
$(1,1,3)$ & $(5,7,23)$ & $-2+12\sigma+\vartheta+8\sigma\vartheta$ & $\frac{11767}{37752}$\\
$(1,1,8)$ & $(8,9,97)$ & $-1+4\sigma\vartheta$ & $\frac{5}{741}$\\
$(1,1,10)$ & $(5,13,110)$ & $-1+4\sigma\vartheta$ & $\frac{400}{9243}$\\
$(1,1,11)$ & $(5,13,111)$ & $1-4\sigma\vartheta$ & $\frac{100}{21093}$\\
$(1,1,12)$ & $(5,13,112)$ & $-1+4\sigma\vartheta$ & $\frac{343}{445}$\\
$(1,2,0)$ & $(5,20,0)$ & $-31-\sigma\vartheta$ & $\frac{57}{1144}$\\
$(1,2,1)$ & $(7,18,7)$ & $\sigma+3\vartheta$ & $\frac{25}{17472}$\\
$(1,2,2)$ & $(6,19,16)$ & $-\sigma+2\vartheta$ & $\frac{75}{4576}$\\
$(1,2,3)$ & $(7,20,29)$ & $-62-3\sigma+31\vartheta-2\sigma\vartheta$ & $\frac{35}{416}$\\
$(1,3,0)$ & $(5,35,4)$ & $-6-2\sigma+4\vartheta-\sigma\vartheta$ & $\frac{665}{858}$\\
$(1,3,1)$ & $(5,34,10)$ & $-2-\vartheta$ & $\frac{11767}{25168}$\\
$(1,3,2)$ & $(7,29,16)$ & $2-3\sigma+\vartheta+2\sigma\vartheta$ & $\frac{35}{416}$\\
$(1,3,3)$ & $(7,26,24)$ & $-\sigma+2\vartheta$ & $\frac{70}{1573}$\\
$(1,4,6)$ & $(6,36,52)$ & $231\sigma-937\vartheta$ & $\frac{49}{585}$\\
$(1,5,5)$ & $(5,44,40)$ & $-\sigma+2\vartheta$ & $\frac{11}{1560}$\\
$(1,5,6)$ & $(6,43,53)$ & $3\sigma-61\vartheta$ & $\frac{5}{3146}$\\
$(1,5,7)$ & $(5,40,68)$ & $-\sigma+2\vartheta$ & $\frac{9}{770}$\\
$(1,6,4)$ & $(6,59,37)$ & $459\sigma-443\vartheta$ & $\frac{49}{585}$\\
$(1,6,5)$ & $(6,58,42)$ & $147\sigma-19\vartheta$ & $\frac{5}{3146}$\\
$(1,6,6)$ & $(5,57,53)$ & $-\sigma+2\vartheta$ & $\frac{416}{1331}$\\
$(1,6,7)$ & $(6,52,68)$ & $-147\sigma+19\vartheta$ & $\frac{125}{363}$\\
$(1,7,5)$ & $(5,81,45)$ & $-\sigma+2\vartheta$ & $\frac{9}{770}$\\
$(1,7,6)$ & $(6,81,59)$ & $-3\sigma+61\vartheta$ & $\frac{125}{363}$\\
$(1,7,7)$ & $(8,69,68)$ & $-\sigma+2\vartheta$ & $\frac{21}{44}$\\
$(1,8,1)$ & $(8,99,9)$ & $1$ & $\frac{5}{741}$\\
$(1,10,1)$ & $(5,110,5)$ & $-1$ & $\frac{400}{9243}$\\
$(1,11,1)$ & $(5,111,5)$ & $1$ & $\frac{100}{21093}$\\
$(1,12,1)$ & $(5,112,5)$ & $-1$ & $\frac{343}{445}$\\
$(2,0,0)$ & $(14,4,0)$ & $-\sigma-3\vartheta$ & $\frac{361}{1560}$\\
$(2,0,1)$ & $(14,2,5)$ & $1-\sigma\vartheta$ & $\frac{133}{572}$\\
$(2,0,2)$ & $(14,2,14)$ & $-1$ & $\frac{91}{380}$\\
$(2,0,3)$ & $(15,0,24)$ & $9+2\sigma+6\vartheta-\sigma\vartheta$ & $\frac{19}{91}$\\
$(2,1,0)$ & $(14,11,0)$ & $-31-\sigma\vartheta$ & $\frac{57}{1144}$\\
$(2,1,1)$ & $(16,9,7)$ & $\sigma+3\vartheta$ & $\frac{25}{17472}$\\
$(2,1,2)$ & $(15,10,16)$ & $-\sigma+2\vartheta$ & $\frac{75}{4576}$\\
$(2,1,3)$ & $(16,11,29)$ & $-62-3\sigma+31\vartheta-2\sigma\vartheta$ & $\frac{35}{416}$\\
$(2,2,0)$ & $(14,20,0)$ & $7-3\sigma\vartheta$ & $\frac{39}{760}$\\
$(2,2,1)$ & $(15,19,7)$ & $\sigma-\vartheta$ & $\frac{75}{4576}$\\
$(2,2,2)$ & $(16,19,17)$ & $-\sigma-3\vartheta$ & $\frac{21}{520}$\\
$(2,2,3)$ & $(14,16,23)$ & $14-9\sigma-7\vartheta-6\sigma\vartheta$ & $\frac{3}{728}$\\
$(2,2,8)$ & $(17,18,97)$ & $7-3\sigma\vartheta$ & $\frac{5}{741}$\\
$(2,2,10)$ & $(14,22,110)$ & $7-3\sigma\vartheta$ & $\frac{400}{9243}$\\
$(2,2,11)$ & $(14,22,111)$ & $7-3\sigma\vartheta$ & $\frac{112908}{175775}$\\
$(2,2,12)$ & $(14,22,112)$ & $-7+3\sigma\vartheta$ & $\frac{45}{623}$\\
$(2,3,0)$ & $(15,34,4)$ & $9-2\sigma-6\vartheta-\sigma\vartheta$ & $\frac{19}{91}$\\
$(2,3,1)$ & $(16,29,7)$ & $2+3\sigma+\vartheta-2\sigma\vartheta$ & $\frac{35}{416}$\\
$(2,3,2)$ & $(14,35,20)$ & $-2-\vartheta$ & $\frac{3}{728}$\\
$(2,3,3)$ & $(16,29,27)$ & $\sigma+3\vartheta$ & $\frac{9}{104}$\\
$(2,4,6)$ & $(17,37,55)$ & $129\sigma+197\vartheta$ & $\frac{11}{65}$\\
$(2,5,5)$ & $(14,44,40)$ & $-\sigma-3\vartheta$ & $\frac{1573}{1176}$\\
$(2,5,6)$ & $(15,44,54)$ & $-9\sigma-5\vartheta$ & $\frac{27}{143}$\\
$(2,5,7)$ & $(14,40,68)$ & $\sigma+3\vartheta$ & $\frac{9}{14}$\\
$(2,6,4)$ & $(17,56,36)$ & $87\sigma+641\vartheta$ & $\frac{11}{65}$\\
$(2,6,5)$ & $(15,57,41)$ & $9\sigma-193\vartheta$ & $\frac{27}{143}$\\
$(2,6,6)$ & $(18,52,52)$ & $-\sigma-3\vartheta$ & $\frac{832}{63525}$\\
$(2,6,7)$ & $(15,52,68)$ & $9\sigma-193\vartheta$ & $\frac{3}{11}$\\
$(2,7,5)$ & $(14,81,45)$ & $\sigma+3\vartheta$ & $\frac{9}{14}$\\
$(2,7,6)$ & $(15,76,54)$ & $9\sigma+5\vartheta$ & $\frac{9}{91}$\\
$(2,7,7)$ & $(17,69,68)$ & $\sigma+3\vartheta$ & $\frac{5}{84}$\\
$(2,8,2)$ & $(17,99,18)$ & $1$ & $\frac{5}{741}$\\
$(2,10,2)$ & $(14,110,14)$ & $-1$ & $\frac{400}{9243}$\\
$(2,11,2)$ & $(14,111,14)$ & $-1$ & $\frac{112908}{175775}$\\
$(2,12,2)$ & $(14,112,14)$ & $1$ & $\frac{45}{623}$\\
$(3,0,1)$ & $(23,4,5)$ & $6-2\sigma-4\vartheta-\sigma\vartheta$ & $\frac{665}{858}$\\
$(3,0,2)$ & $(24,4,15)$ & $9+2\sigma-6\vartheta+\sigma\vartheta$ & $\frac{19}{91}$\\
$(3,0,3)$ & $(24,2,24)$ & $1$ & $\frac{80}{1729}$\\
$(3,1,0)$ & $(23,13,0)$ & $-6-2\sigma+4\vartheta-\sigma\vartheta$ & $\frac{665}{858}$\\
$(3,1,1)$ & $(23,11,5)$ & $-2-\vartheta$ & $\frac{11767}{37752}$\\
$(3,1,2)$ & $(29,7,16)$ & $2-3\sigma+\vartheta+2\sigma\vartheta$ & $\frac{35}{416}$\\
$(3,1,3)$ & $(24,9,24)$ & $\sigma-2\vartheta$ & $\frac{35}{1716}$\\
$(3,2,0)$ & $(24,21,0)$ & $9-2\sigma-6\vartheta-\sigma\vartheta$ & $\frac{19}{91}$\\
$(3,2,1)$ & $(29,16,7)$ & $2+3\sigma+\vartheta-2\sigma\vartheta$ & $\frac{35}{416}$\\
$(3,2,2)$ & $(23,20,14)$ & $-2-\vartheta$ & $\frac{3}{728}$\\
$(3,2,3)$ & $(27,20,29)$ & $\sigma+3\vartheta$ & $\frac{9}{104}$\\
$(3,3,0)$ & $(24,34,2)$ & $-1-4\vartheta$ & $\frac{80}{1729}$\\
$(3,3,1)$ & $(24,32,7)$ & $-24+\sigma+2\vartheta+4\sigma\vartheta$ & $\frac{70}{1573}$\\
$(3,3,2)$ & $(27,29,16)$ & $36+\sigma-3\vartheta+4\sigma\vartheta$ & $\frac{9}{104}$\\
$(3,3,3)$ & $(23,30,24)$ & $2+\vartheta$ & $\frac{9}{143}$\\
$(3,3,8)$ & $(25,33,98)$ & $1+4\vartheta$ & $\frac{2}{247}$\\
$(3,3,10)$ & $(23,35,110)$ & $-1-4\vartheta$ & $\frac{196}{1027}$\\
$(3,3,11)$ & $(23,35,111)$ & $1+4\vartheta$ & $\frac{147}{7031}$\\
$(3,3,12)$ & $(23,35,112)$ & $-1-4\vartheta$ & $\frac{5}{623}$\\
$(3,4,7)$ & $(23,36,68)$ & $-32+3\sigma-16\vartheta-2\sigma\vartheta$ & $\frac{1}{2}$\\
$(3,5,6)$ & $(24,45,53)$ & $-162-33\sigma-81\vartheta+22\sigma\vartheta$ & $\frac{128}{429}$\\
$(3,5,7)$ & $(23,42,68)$ & $-2-\vartheta$ & $\frac{1}{7}$\\
$(3,6,5)$ & $(24,58,40)$ & $-162+33\sigma-81\vartheta-22\sigma\vartheta$ & $\frac{128}{429}$\\
$(3,6,6)$ & $(25,57,53)$ & $2+\vartheta$ & $\frac{5776}{51909}$\\
$(3,6,7)$ & $(23,55,68)$ & $-162+33\sigma-81\vartheta-22\sigma\vartheta$ & $\frac{24}{77}$\\
$(3,7,4)$ & $(23,81,37)$ & $-32-3\sigma-16\vartheta+2\sigma\vartheta$ & $\frac{1}{2}$\\
$(3,7,5)$ & $(23,81,43)$ & $-2-\vartheta$ & $\frac{1}{7}$\\
$(3,7,6)$ & $(23,81,56)$ & $162+33\sigma+81\vartheta-22\sigma\vartheta$ & $\frac{24}{77}$\\
$(3,7,7)$ & $(23,74,68)$ & $2+\vartheta$ & $\frac{1}{22}$\\
$(3,8,3)$ & $(25,98,25)$ & $1$ & $\frac{2}{247}$\\
$(3,10,3)$ & $(23,110,23)$ & $-1$ & $\frac{196}{1027}$\\
$(3,11,3)$ & $(23,111,23)$ & $1$ & $\frac{147}{7031}$\\
$(3,12,3)$ & $(23,112,23)$ & $-1$ & $\frac{5}{623}$\\
$(4,0,5)$ & $(36,0,40)$ & $-16-\sigma\vartheta$ & $\frac{19}{78}$\\
$(4,1,6)$ & $(36,6,52)$ & $231\sigma-937\vartheta$ & $\frac{49}{585}$\\
$(4,2,6)$ & $(36,19,56)$ & $-129\sigma-197\vartheta$ & $\frac{11}{65}$\\
$(4,3,7)$ & $(36,23,68)$ & $-32+3\sigma-16\vartheta-2\sigma\vartheta$ & $\frac{1}{2}$\\
$(4,4,9)$ & $(36,36,101)$ & $-217-32\sigma\vartheta$ & $\frac{1}{39}$\\
$(4,4,10)$ & $(36,39,110)$ & $-217-32\sigma\vartheta$ & $\frac{361}{4108}$\\
$(4,4,11)$ & $(36,39,111)$ & $217+32\sigma\vartheta$ & $\frac{48}{175775}$\\
$(4,4,12)$ & $(36,39,112)$ & $217+32\sigma\vartheta$ & $\frac{567}{445}$\\
$(4,5,0)$ & $(36,51,4)$ & $16+\sigma\vartheta$ & $\frac{19}{78}$\\
$(4,6,1)$ & $(36,62,7)$ & $-27\sigma+203\vartheta$ & $\frac{98}{195}$\\
$(4,6,2)$ & $(36,63,17)$ & $-159\sigma-37\vartheta$ & $\frac{11}{65}$\\
$(4,7,3)$ & $(36,95,35)$ & $32+3\sigma-16\vartheta+2\sigma\vartheta$ & $\frac{1}{2}$\\
$(4,9,4)$ & $(36,106,37)$ & $1$ & $\frac{1}{78}$\\
$(4,10,4)$ & $(36,110,36)$ & $1$ & $\frac{361}{4108}$\\
$(4,11,4)$ & $(36,111,36)$ & $-1$ & $\frac{48}{175775}$\\
$(4,12,4)$ & $(36,112,36)$ & $-1$ & $\frac{567}{445}$\\
$(5,0,4)$ & $(40,4,36)$ & $-16+\sigma\vartheta$ & $\frac{19}{78}$\\
$(5,0,5)$ & $(41,2,41)$ & $1$ & $\frac{605}{338884}$\\
$(5,0,6)$ & $(41,2,54)$ & $-81+11\sigma\vartheta$ & $\frac{836}{4459}$\\
$(5,1,5)$ & $(40,10,41)$ & $-\sigma+2\vartheta$ & $\frac{11}{5460}$\\
$(5,1,6)$ & $(41,8,53)$ & $-3\sigma+61\vartheta$ & $\frac{65}{847}$\\
$(5,1,7)$ & $(40,5,68)$ & $-\sigma+2\vartheta$ & $\frac{9}{770}$\\
$(5,2,5)$ & $(40,19,41)$ & $-\sigma-3\vartheta$ & $\frac{1573}{4116}$\\
$(5,2,6)$ & $(41,21,57)$ & $9\sigma+5\vartheta$ & $\frac{27}{143}$\\
$(5,2,7)$ & $(40,14,68)$ & $\sigma+3\vartheta$ & $\frac{9}{14}$\\
$(5,3,6)$ & $(40,30,54)$ & $162+33\sigma+81\vartheta-22\sigma\vartheta$ & $\frac{64}{4719}$\\
$(5,3,7)$ & $(42,23,68)$ & $-2-\vartheta$ & $\frac{1}{7}$\\
$(5,4,0)$ & $(40,39,0)$ & $16+\sigma\vartheta$ & $\frac{19}{78}$\\
$(5,5,0)$ & $(41,49,1)$ & $-1$ & $\frac{1815}{84721}$\\
$(5,5,1)$ & $(40,47,5)$ & $\sigma+2\vartheta$ & $\frac{11}{1560}$\\
$(5,5,2)$ & $(40,47,14)$ & $\sigma-3\vartheta$ & $\frac{1573}{1176}$\\
$(5,5,8)$ & $(41,48,96)$ & $1$ & $\frac{6}{247}$\\
$(5,5,9)$ & $(40,44,101)$ & $-1$ & $\frac{5}{39}$\\
$(5,5,10)$ & $(40,51,110)$ & $-1$ & $\frac{1681}{36972}$\\
$(5,5,11)$ & $(40,51,111)$ & $1$ & $\frac{29403}{28124}$\\
$(5,5,12)$ & $(40,51,112)$ & $-1$ & $\frac{841}{3115}$\\
$(5,6,0)$ & $(41,65,2)$ & $81+11\sigma\vartheta$ & $\frac{836}{4459}$\\
$(5,6,1)$ & $(41,66,10)$ & $3\sigma+61\vartheta$ & $\frac{65}{847}$\\
$(5,6,2)$ & $(41,62,15)$ & $-9\sigma+5\vartheta$ & $\frac{27}{143}$\\
$(5,6,3)$ & $(40,61,24)$ & $162+33\sigma-81\vartheta+22\sigma\vartheta$ & $\frac{128}{429}$\\
$(5,7,1)$ & $(40,93,11)$ & $\sigma+2\vartheta$ & $\frac{3}{715}$\\
$(5,7,2)$ & $(40,93,20)$ & $-\sigma+3\vartheta$ & $\frac{3}{13}$\\
$(5,7,3)$ & $(42,95,35)$ & $2-\vartheta$ & $\frac{1}{7}$\\
$(5,8,5)$ & $(41,99,42)$ & $1$ & $\frac{8}{741}$\\
$(5,9,5)$ & $(40,106,41)$ & $1$ & $\frac{5}{78}$\\
$(5,10,5)$ & $(40,110,40)$ & $1$ & $\frac{1681}{36972}$\\
$(5,11,5)$ & $(40,111,40)$ & $-1$ & $\frac{29403}{28124}$\\
$(5,12,5)$ & $(40,112,40)$ & $1$ & $\frac{841}{3115}$\\
$(6,0,5)$ & $(54,2,41)$ & $-81-11\sigma\vartheta$ & $\frac{836}{4459}$\\
$(6,0,6)$ & $(52,2,52)$ & $1$ & $\frac{208}{5985}$\\
$(6,1,4)$ & $(52,12,36)$ & $-459\sigma+443\vartheta$ & $\frac{49}{585}$\\
$(6,1,5)$ & $(53,10,41)$ & $147\sigma-19\vartheta$ & $\frac{65}{847}$\\
$(6,1,6)$ & $(52,10,53)$ & $-\sigma+2\vartheta$ & $\frac{832}{3993}$\\
$(6,1,7)$ & $(52,6,68)$ & $-147\sigma+19\vartheta$ & $\frac{125}{363}$\\
$(6,2,4)$ & $(55,19,37)$ & $-87\sigma-641\vartheta$ & $\frac{11}{65}$\\
$(6,2,5)$ & $(54,21,44)$ & $-9\sigma+193\vartheta$ & $\frac{27}{143}$\\
$(6,2,6)$ & $(52,18,52)$ & $-\sigma-3\vartheta$ & $\frac{832}{63525}$\\
$(6,2,7)$ & $(52,15,68)$ & $9\sigma-193\vartheta$ & $\frac{3}{11}$\\
$(6,3,5)$ & $(53,34,45)$ & $162-33\sigma+81\vartheta+22\sigma\vartheta$ & $\frac{128}{429}$\\
$(6,3,6)$ & $(53,33,57)$ & $2+\vartheta$ & $\frac{5776}{51909}$\\
$(6,3,7)$ & $(52,26,68)$ & $162-33\sigma+81\vartheta+22\sigma\vartheta$ & $\frac{6}{121}$\\
$(6,4,1)$ & $(52,39,6)$ & $27\sigma-203\vartheta$ & $\frac{49}{585}$\\
$(6,4,2)$ & $(55,38,17)$ & $159\sigma+37\vartheta$ & $\frac{11}{65}$\\
$(6,5,0)$ & $(54,50,2)$ & $-81-11\sigma\vartheta$ & $\frac{836}{4459}$\\
$(6,5,1)$ & $(53,48,6)$ & $3\sigma+61\vartheta$ & $\frac{5}{3146}$\\
$(6,5,2)$ & $(54,47,15)$ & $9\sigma-5\vartheta$ & $\frac{27}{143}$\\
$(6,5,3)$ & $(53,46,24)$ & $-162-33\sigma+81\vartheta-22\sigma\vartheta$ & $\frac{128}{429}$\\
$(6,6,0)$ & $(52,67,2)$ & $-307-297\sigma\vartheta$ & $\frac{208}{5985}$\\
$(6,6,1)$ & $(52,65,7)$ & $-59\sigma+179\vartheta$ & $\frac{416}{1331}$\\
$(6,6,2)$ & $(52,67,18)$ & $149\sigma+147\vartheta$ & $\frac{832}{63525}$\\
$(6,6,3)$ & $(53,62,25)$ & $614+891\sigma-307\vartheta+594\sigma\vartheta$ & $\frac{5776}{51909}$\\
$(6,6,8)$ & $(52,65,96)$ & $-307-297\sigma\vartheta$ & $\frac{7}{247}$\\
$(6,6,9)$ & $(52,58,101)$ & $-307-297\sigma\vartheta$ & $\frac{7}{39}$\\
$(6,6,10)$ & $(52,67,110)$ & $-307-297\sigma\vartheta$ & $\frac{81}{4108}$\\
$(6,6,11)$ & $(52,67,111)$ & $307+297\sigma\vartheta$ & $\frac{29584}{527325}$\\
$(6,6,12)$ & $(52,67,112)$ & $307+297\sigma\vartheta$ & $\frac{169}{3115}$\\
$(6,7,1)$ & $(52,95,12)$ & $-3\sigma-61\vartheta$ & $\frac{125}{363}$\\
$(6,7,2)$ & $(52,95,21)$ & $-9\sigma+5\vartheta$ & $\frac{3}{11}$\\
$(6,7,3)$ & $(52,95,32)$ & $162+33\sigma-81\vartheta+22\sigma\vartheta$ & $\frac{6}{121}$\\
$(6,8,6)$ & $(52,100,54)$ & $1$ & $\frac{7}{247}$\\
$(6,9,6)$ & $(52,106,53)$ & $1$ & $\frac{7}{78}$\\
$(6,10,6)$ & $(52,110,52)$ & $1$ & $\frac{81}{4108}$\\
$(6,11,6)$ & $(52,111,52)$ & $-1$ & $\frac{29584}{527325}$\\
$(6,12,6)$ & $(52,112,52)$ & $-1$ & $\frac{169}{3115}$\\
$(7,0,7)$ & $(68,4,70)$ & $-1$ & $\frac{5}{133}$\\
$(7,1,5)$ & $(68,13,40)$ & $-\sigma+2\vartheta$ & $\frac{9}{770}$\\
$(7,1,6)$ & $(68,12,52)$ & $3\sigma-61\vartheta$ & $\frac{125}{363}$\\
$(7,1,7)$ & $(68,10,69)$ & $\sigma-2\vartheta$ & $\frac{21}{44}$\\
$(7,2,5)$ & $(68,22,40)$ & $\sigma+3\vartheta$ & $\frac{9}{14}$\\
$(7,2,6)$ & $(68,21,52)$ & $9\sigma+5\vartheta$ & $\frac{3}{11}$\\
$(7,2,7)$ & $(68,19,69)$ & $-\sigma-3\vartheta$ & $\frac{5}{84}$\\
$(7,3,4)$ & $(68,35,36)$ & $-32-3\sigma-16\vartheta+2\sigma\vartheta$ & $\frac{1}{2}$\\
$(7,3,5)$ & $(68,35,42)$ & $-2-\vartheta$ & $\frac{1}{7}$\\
$(7,3,6)$ & $(68,32,52)$ & $-162-33\sigma-81\vartheta+22\sigma\vartheta$ & $\frac{6}{121}$\\
$(7,3,7)$ & $(68,30,69)$ & $2+\vartheta$ & $\frac{3}{11}$\\
$(7,4,3)$ & $(68,39,23)$ & $32+3\sigma-16\vartheta+2\sigma\vartheta$ & $\frac{1}{2}$\\
$(7,5,1)$ & $(68,51,5)$ & $\sigma+2\vartheta$ & $\frac{9}{770}$\\
$(7,5,2)$ & $(68,51,14)$ & $-\sigma+3\vartheta$ & $\frac{9}{14}$\\
$(7,5,3)$ & $(68,49,23)$ & $2-\vartheta$ & $\frac{1}{7}$\\
$(7,6,1)$ & $(68,67,6)$ & $3\sigma+61\vartheta$ & $\frac{125}{363}$\\
$(7,6,2)$ & $(68,67,15)$ & $9\sigma-5\vartheta$ & $\frac{3}{11}$\\
$(7,6,3)$ & $(68,64,23)$ & $-162-33\sigma+81\vartheta-22\sigma\vartheta$ & $\frac{24}{77}$\\
$(7,7,0)$ & $(68,93,0)$ & $1$ & $\frac{5}{133}$\\
$(7,7,1)$ & $(68,94,8)$ & $-\sigma-2\vartheta$ & $\frac{21}{44}$\\
$(7,7,2)$ & $(68,94,17)$ & $\sigma-3\vartheta$ & $\frac{5}{84}$\\
$(7,7,3)$ & $(68,89,23)$ & $-2+\vartheta$ & $\frac{1}{22}$\\
$(7,7,8)$ & $(68,93,96)$ & $-1$ & $\frac{2}{19}$\\
$(7,7,9)$ & $(68,80,101)$ & $-1$ & $\frac{1}{3}$\\
$(7,7,10)$ & $(68,95,110)$ & $1$ & $\frac{13}{316}$\\
$(7,7,11)$ & $(68,95,111)$ & $-1$ & $\frac{1587}{7031}$\\
$(7,7,12)$ & $(68,95,112)$ & $-1$ & $\frac{81}{3115}$\\
$(7,8,7)$ & $(68,100,70)$ & $1$ & $\frac{2}{19}$\\
$(7,9,7)$ & $(68,106,69)$ & $1$ & $\frac{1}{6}$\\
$(7,10,7)$ & $(68,110,68)$ & $-1$ & $\frac{13}{316}$\\
$(7,11,7)$ & $(68,111,68)$ & $1$ & $\frac{1587}{7031}$\\
$(7,12,7)$ & $(68,112,68)$ & $1$ & $\frac{81}{3115}$\\
$(8,0,0)$ & $(96,2,0)$ & $1$ & $\frac{2}{247}$\\
$(8,1,1)$ & $(97,9,8)$ & $-1$ & $\frac{5}{741}$\\
$(8,2,2)$ & $(97,18,17)$ & $-1$ & $\frac{5}{741}$\\
$(8,3,3)$ & $(97,29,28)$ & $-1$ & $\frac{7}{494}$\\
$(8,5,5)$ & $(96,43,41)$ & $1$ & $\frac{6}{247}$\\
$(8,6,6)$ & $(96,54,52)$ & $1$ & $\frac{7}{247}$\\
$(8,7,7)$ & $(96,70,68)$ & $1$ & $\frac{2}{19}$\\
$(8,8,8)$ & $(96,99,97)$ & $-1$ & $\frac{1225}{741}$\\
$(8,8,10)$ & $(96,100,110)$ & $-1$ & $\frac{6724}{9243}$\\
$(8,8,11)$ & $(96,100,111)$ & $1$ & $\frac{4800}{7031}$\\
$(8,8,12)$ & $(96,100,112)$ & $1$ & $\frac{35}{89}$\\
$(8,9,9)$ & $(96,103,101)$ & $-1$ & $\frac{76}{39}$\\
$(8,10,8)$ & $(96,110,96)$ & $-1$ & $\frac{6724}{9243}$\\
$(8,11,8)$ & $(96,111,96)$ & $1$ & $\frac{4800}{7031}$\\
$(8,12,8)$ & $(96,112,96)$ & $1$ & $\frac{35}{89}$\\
$(9,4,4)$ & $(101,39,36)$ & $1$ & $\frac{1}{39}$\\
$(9,5,5)$ & $(101,47,40)$ & $1$ & $\frac{5}{39}$\\
$(9,6,6)$ & $(101,61,52)$ & $1$ & $\frac{7}{39}$\\
$(9,7,7)$ & $(101,83,68)$ & $1$ & $\frac{1}{3}$\\
$(9,8,9)$ & $(101,99,102)$ & $-1$ & $\frac{38}{39}$\\
$(9,9,8)$ & $(101,107,96)$ & $-1$ & $\frac{76}{39}$\\
$(9,9,9)$ & $(101,105,101)$ & $1$ & $\frac{4}{39}$\\
$(9,9,10)$ & $(101,109,110)$ & $1$ & $\frac{5776}{9243}$\\
$(9,9,11)$ & $(101,109,111)$ & $1$ & $\frac{4800}{7031}$\\
$(9,9,12)$ & $(101,109,112)$ & $1$ & $\frac{35}{89}$\\
$(9,10,9)$ & $(101,110,101)$ & $1$ & $\frac{5776}{9243}$\\
$(9,11,9)$ & $(101,111,101)$ & $1$ & $\frac{4800}{7031}$\\
$(9,12,9)$ & $(101,112,101)$ & $1$ & $\frac{35}{89}$\\
$(10,0,0)$ & $(110,0,0)$ & $-1$ & $\frac{4}{1027}$\\
$(10,1,1)$ & $(110,5,5)$ & $-1$ & $\frac{400}{9243}$\\
$(10,2,2)$ & $(110,14,14)$ & $-1$ & $\frac{400}{9243}$\\
$(10,3,3)$ & $(110,23,23)$ & $-1$ & $\frac{196}{1027}$\\
$(10,4,4)$ & $(110,36,36)$ & $1$ & $\frac{361}{4108}$\\
$(10,5,5)$ & $(110,40,40)$ & $1$ & $\frac{1681}{36972}$\\
$(10,6,6)$ & $(110,52,52)$ & $1$ & $\frac{81}{4108}$\\
$(10,7,7)$ & $(110,68,68)$ & $-1$ & $\frac{13}{316}$\\
$(10,8,8)$ & $(110,96,96)$ & $-1$ & $\frac{6724}{9243}$\\
$(10,9,9)$ & $(110,101,101)$ & $1$ & $\frac{5776}{9243}$\\
$(10,10,10)$ & $(110,110,110)$ & $1$ & $\frac{23104}{9243}$\\
$(10,10,11)$ & $(110,110,111)$ & $1$ & $\frac{4800}{7031}$\\
$(10,10,12)$ & $(110,110,112)$ & $1$ & $\frac{35}{89}$\\
$(10,11,10)$ & $(110,111,110)$ & $1$ & $\frac{4800}{7031}$\\
$(10,12,10)$ & $(110,112,110)$ & $1$ & $\frac{35}{89}$\\
$(11,0,0)$ & $(111,0,0)$ & $-1$ & $\frac{154587}{703100}$\\
$(11,1,1)$ & $(111,5,5)$ & $1$ & $\frac{100}{21093}$\\
$(11,2,2)$ & $(111,14,14)$ & $-1$ & $\frac{112908}{175775}$\\
$(11,3,3)$ & $(111,23,23)$ & $1$ & $\frac{147}{7031}$\\
$(11,4,4)$ & $(111,36,36)$ & $-1$ & $\frac{48}{175775}$\\
$(11,5,5)$ & $(111,40,40)$ & $-1$ & $\frac{29403}{28124}$\\
$(11,6,6)$ & $(111,52,52)$ & $-1$ & $\frac{29584}{527325}$\\
$(11,7,7)$ & $(111,68,68)$ & $1$ & $\frac{1587}{7031}$\\
$(11,8,8)$ & $(111,96,96)$ & $1$ & $\frac{4800}{7031}$\\
$(11,9,9)$ & $(111,101,101)$ & $1$ & $\frac{4800}{7031}$\\
$(11,10,10)$ & $(111,110,110)$ & $1$ & $\frac{4800}{7031}$\\
$(11,11,11)$ & $(111,111,111)$ & $-1$ & $\frac{1560001}{2109300}$\\
$(11,11,12)$ & $(111,111,112)$ & $1$ & $\frac{35}{89}$\\
$(11,12,11)$ & $(111,112,111)$ & $1$ & $\frac{35}{89}$\\
$(12,0,0)$ & $(112,0,0)$ & $-1$ & $\frac{1681}{3115}$\\
$(12,1,1)$ & $(112,5,5)$ & $-1$ & $\frac{343}{445}$\\
$(12,2,2)$ & $(112,14,14)$ & $1$ & $\frac{45}{623}$\\
$(12,3,3)$ & $(112,23,23)$ & $-1$ & $\frac{5}{623}$\\
$(12,4,4)$ & $(112,36,36)$ & $-1$ & $\frac{567}{445}$\\
$(12,5,5)$ & $(112,40,40)$ & $1$ & $\frac{841}{3115}$\\
$(12,6,6)$ & $(112,52,52)$ & $-1$ & $\frac{169}{3115}$\\
$(12,7,7)$ & $(112,68,68)$ & $1$ & $\frac{81}{3115}$\\
$(12,8,8)$ & $(112,96,96)$ & $1$ & $\frac{35}{89}$\\
$(12,9,9)$ & $(112,101,101)$ & $1$ & $\frac{35}{89}$\\
$(12,10,10)$ & $(112,110,110)$ & $1$ & $\frac{35}{89}$\\
$(12,11,11)$ & $(112,111,111)$ & $1$ & $\frac{35}{89}$\\
$(12,12,12)$ & $(112,112,112)$ & $-1$ & $\frac{2916}{3115}$\\
\end{longtable}
\endgroup


\subsection{Stress projections and an independent four-point cross-check}

For a unit primary of internal stress weight $h_i$, the squared
coupling to the unit stress $T_i/\sqrt{c_i/2}$ is $2h_i^2/c_i$.
The complete saved list for the eight original current multiplets
is printed below. It resolves the two distinct $(1,9)$ multiplets
without relying on their names. Each multiplet's four $h_i$ values
sum to two.

\begingroup\small
\setlength{\tabcolsep}{5pt}
\begin{longtable}{rlrrr}
\caption{Resolved stress couplings for the eight original highest multiplets. These are intrinsic-stress projections, before changing to the total stress and the three orthogonal stress-difference primaries.}\label{expphysics:stress-ward}\\
\toprule
$m$ & stress & $c_i$ & $h_i$ & $2h_i^2/c_i$\\
\midrule
\endfirsthead
\multicolumn{5}{l}{\small\itshape Continued from preceding page.}\\
\toprule
$m$ & stress & $c_i$ & $h_i$ & $2h_i^2/c_i$\\
\midrule
\endhead
\midrule
\multicolumn{5}{r}{\small\itshape Continued on next page.}\\
\endfoot
\bottomrule
\endlastfoot
0 & $T_L$ & $1$ & $0$ & $0$\\
0 & $T_R$ & $\frac{117}{41}$ & $\frac{6}{41}$ & $\frac{8}{533}$\\
0 & $T_e$ & $\frac{2400}{1271}$ & $\frac{798}{1271}$ & $\frac{53067}{127100}$\\
0 & $T_f$ & $\frac{70}{31}$ & $\frac{38}{31}$ & $\frac{1444}{1085}$\\
1 & $T_L$ & $1$ & $0$ & $0$\\
1 & $T_R$ & $\frac{117}{41}$ & $\frac{20}{41}$ & $\frac{800}{4797}$\\
1 & $T_e$ & $\frac{2400}{1271}$ & $\frac{200}{1271}$ & $\frac{100}{3813}$\\
1 & $T_f$ & $\frac{70}{31}$ & $\frac{42}{31}$ & $\frac{252}{155}$\\
2 & $T_L$ & $1$ & $0$ & $0$\\
2 & $T_R$ & $\frac{117}{41}$ & $\frac{20}{41}$ & $\frac{800}{4797}$\\
2 & $T_e$ & $\frac{2400}{1271}$ & $\frac{1512}{1271}$ & $\frac{47628}{31775}$\\
2 & $T_f$ & $\frac{70}{31}$ & $\frac{10}{31}$ & $\frac{20}{217}$\\
3 & $T_L$ & $1$ & $0$ & $0$\\
3 & $T_R$ & $\frac{117}{41}$ & $\frac{42}{41}$ & $\frac{392}{533}$\\
3 & $T_e$ & $\frac{2400}{1271}$ & $\frac{420}{1271}$ & $\frac{147}{1271}$\\
3 & $T_f$ & $\frac{70}{31}$ & $\frac{20}{31}$ & $\frac{80}{217}$\\
4 & $T_L$ & $1$ & $\frac{1}{4}$ & $\frac{1}{8}$\\
4 & $T_R$ & $\frac{117}{41}$ & $\frac{3}{164}$ & $\frac{1}{4264}$\\
4 & $T_e$ & $\frac{2400}{1271}$ & $\frac{192}{1271}$ & $\frac{768}{31775}$\\
4 & $T_f$ & $\frac{70}{31}$ & $\frac{49}{31}$ & $\frac{343}{155}$\\
5 & $T_L$ & $1$ & $\frac{1}{4}$ & $\frac{1}{8}$\\
5 & $T_R$ & $\frac{117}{41}$ & $\frac{35}{164}$ & $\frac{1225}{38376}$\\
5 & $T_e$ & $\frac{2400}{1271}$ & $\frac{1830}{1271}$ & $\frac{11163}{5084}$\\
5 & $T_f$ & $\frac{70}{31}$ & $\frac{3}{31}$ & $\frac{9}{1085}$\\
6 & $T_L$ & $1$ & $\frac{1}{4}$ & $\frac{1}{8}$\\
6 & $T_R$ & $\frac{117}{41}$ & $\frac{63}{164}$ & $\frac{441}{4264}$\\
6 & $T_e$ & $\frac{2400}{1271}$ & $\frac{752}{1271}$ & $\frac{35344}{95325}$\\
6 & $T_f$ & $\frac{70}{31}$ & $\frac{24}{31}$ & $\frac{576}{1085}$\\
7 & $T_L$ & $1$ & $\frac{1}{4}$ & $\frac{1}{8}$\\
7 & $T_R$ & $\frac{117}{41}$ & $\frac{195}{164}$ & $\frac{325}{328}$\\
7 & $T_e$ & $\frac{2400}{1271}$ & $\frac{180}{1271}$ & $\frac{27}{1271}$\\
7 & $T_f$ & $\frac{70}{31}$ & $\frac{13}{31}$ & $\frac{169}{1085}$\\
\end{longtable}
\endgroup


The extremal pair $(68,95)$ is
$(v_{7,0,0},v_{7,1,13})$. The normalized second component is the
chiral adjoint of $\widehat u$ up to a phase, because that extremal
Cartan weight occurs only once. There are exactly nine nonzero
primary targets. Printing all nine makes the aggregate
four-point check a short rational addition.

\begingroup\small
\setlength{\tabcolsep}{5pt}
\begin{longtable}{rrr}
\caption{All nine nonzero primary projections of the extremal $u$ pair. Their sum is $11/3$; targets 110, 111, 112 sum to $41/140$.}\label{expphysics:u-pair}\\
\toprule
target $k$ & $(m,p,q)$ & $|C_{68,95}^{\ k}|^2$\\
\midrule
\endfirsthead
\multicolumn{3}{l}{\small\itshape Continued from preceding page.}\\
\toprule
target $k$ & $(m,p,q)$ & $|C_{68,95}^{\ k}|^2$\\
\midrule
\endhead
\midrule
\multicolumn{3}{r}{\small\itshape Continued on next page.}\\
\endfoot
\bottomrule
\endlastfoot
2 & $(0,0,2)$ & $\frac{65}{133}$\\
9 & $(1,0,4)$ & $\frac{273}{880}$\\
18 & $(2,0,4)$ & $\frac{13}{336}$\\
29 & $(3,0,6)$ & $\frac{13}{154}$\\
98 & $(8,0,2)$ & $\frac{26}{19}$\\
105 & $(9,1,1)$ & $\frac{13}{12}$\\
110 & $(10,0,0)$ & $\frac{13}{316}$\\
111 & $(11,0,0)$ & $\frac{1587}{7031}$\\
112 & $(12,0,0)$ & $\frac{81}{3115}$\\
\end{longtable}
\endgroup


Their sum is $11/3$. The three stress-difference targets sum to
$41/140$; adding the total stress contribution 1 gives $181/140$
for the four commuting stresses. The six remaining primary targets
sum to $1417/420$. Thus the total global weight-two coefficient is
\[
 \frac{14}{3}
 =1+\frac{11}{3}
 =\frac{181}{140}+\frac{1417}{420}.
\]
The two decompositions resolve different orthogonal subspaces.
Neither interpretation counts the derivative of a weight-one
current as a new global quasiprimary.

\clearpage
\section{Correlator verification tables and exact channel expansions}
\label{expphysics:correlator-details}

This section prints the channel data of
\artifact{c8_physics_20260927/correlators.json}, including all
coefficients through the saved cutoff 32. Rational-function
identities are exact at every cross-ratio where defined. A finite
global-block table records coefficients only through its stated
cutoff. Those two statements should be kept distinct when checking
crossing and positivity.

\subsection{A Gaussian calculation with every surviving term visible}

Let the three lattice momenta of the defining weight-two field be
\begin{align*}
 a_0&=(0,0,0,0,0,-1,1,1),\\
 a_1&=(0,0,0,1,1,-1,0,0),\\
 a_2&=(1,1,0,0,0,-1,0,0).
\end{align*}
They have Gram matrix with diagonal 4 and off-diagonal 2, and
\[
 u=e^{a_0}+\frac{\sqrt{13}+3\sqrt{-3}}{20}e^{a_1}
            +\frac{\sqrt{13}-3\sqrt{-3}}{20}e^{a_2}.
\]
The three squared exponential norms in the transported form are
$(4,40,40)$. The coefficient absolute squares are $(1,1/10,1/10)$,
so each contributes four to the norm $\|u\|^2=12$. The normalized
field is $\widehat u=u/\sqrt{12}$. An adjoint at a chiral insertion
is still a holomorphic field; it is not an antiholomorphic insertion.

Number the four insertion points $(\infty,1,x,0)$. A letter $p$
means $\widehat u$ and a letter $m$ its chiral adjoint. Inserting
the three exponential terms at each position gives 81 choices per
ordering. Momentum conservation is imposed before assigning a
nonzero correlator. Exactly fifteen choices survive in each of the
six two-$p$, two-$m$ orderings. Although the cocycle signs vary,
the complete normalized coefficient of each surviving monomial is
$1/9$. The cocycle sign in the next table is one factor in this
calculation, not the sign of the final correlator contribution.

\begingroup\small
\setlength{\tabcolsep}{5pt}
\begin{longtable}{lrrrr}
\caption{Every surviving Gaussian term in the two independent reflection-positive $u$ orderings. Choices are numbered 0,1,2 in the momentum order of the text; each complete normalized coefficient is $1/9$, including its adjoint and cocycle factors.}\label{expphysics:gaussian-terms}\\
\toprule
ordering & momentum choices & cocycle sign & $x$ exponent & $(1-x)$ exponent\\
\midrule
\endfirsthead
\multicolumn{5}{l}{\small\itshape Continued from preceding page.}\\
\toprule
ordering & momentum choices & cocycle sign & $x$ exponent & $(1-x)$ exponent\\
\midrule
\endhead
\midrule
\multicolumn{5}{r}{\small\itshape Continued on next page.}\\
\endfoot
\bottomrule
\endlastfoot
\texttt{mpmp} & $(0,0,0,0)$ & 1 & -4 & -4\\
\texttt{mpmp} & $(0,0,1,1)$ & 1 & -4 & -2\\
\texttt{mpmp} & $(0,0,2,2)$ & 1 & -4 & -2\\
\texttt{mpmp} & $(0,1,1,0)$ & 1 & -2 & -4\\
\texttt{mpmp} & $(0,2,2,0)$ & 1 & -2 & -4\\
\texttt{mpmp} & $(1,0,0,1)$ & 1 & -2 & -4\\
\texttt{mpmp} & $(1,1,0,0)$ & 1 & -4 & -2\\
\texttt{mpmp} & $(1,1,1,1)$ & 1 & -4 & -4\\
\texttt{mpmp} & $(1,1,2,2)$ & 1 & -4 & -2\\
\texttt{mpmp} & $(1,2,2,1)$ & 1 & -2 & -4\\
\texttt{mpmp} & $(2,0,0,2)$ & 1 & -2 & -4\\
\texttt{mpmp} & $(2,1,1,2)$ & 1 & -2 & -4\\
\texttt{mpmp} & $(2,2,0,0)$ & 1 & -4 & -2\\
\texttt{mpmp} & $(2,2,1,1)$ & 1 & -4 & -2\\
\texttt{mpmp} & $(2,2,2,2)$ & 1 & -4 & -4\\
\texttt{mmpp} & $(0,0,0,0)$ & 1 & 4 & -4\\
\texttt{mmpp} & $(0,1,0,1)$ & 1 & 2 & -2\\
\texttt{mmpp} & $(0,1,1,0)$ & 1 & 2 & -4\\
\texttt{mmpp} & $(0,2,0,2)$ & 1 & 2 & -2\\
\texttt{mmpp} & $(0,2,2,0)$ & 1 & 2 & -4\\
\texttt{mmpp} & $(1,0,0,1)$ & 1 & 2 & -4\\
\texttt{mmpp} & $(1,0,1,0)$ & 1 & 2 & -2\\
\texttt{mmpp} & $(1,1,1,1)$ & 1 & 4 & -4\\
\texttt{mmpp} & $(1,2,1,2)$ & 1 & 2 & -2\\
\texttt{mmpp} & $(1,2,2,1)$ & 1 & 2 & -4\\
\texttt{mmpp} & $(2,0,0,2)$ & 1 & 2 & -4\\
\texttt{mmpp} & $(2,0,2,0)$ & 1 & 2 & -2\\
\texttt{mmpp} & $(2,1,1,2)$ & 1 & 2 & -4\\
\texttt{mmpp} & $(2,1,2,1)$ & 1 & 2 & -2\\
\texttt{mmpp} & $(2,2,2,2)$ & 1 & 4 & -4\\
\end{longtable}
\endgroup


Summing the monomials gives
\begin{align}
 G_0(x)&=G_{mpmp}(x)=G_{pmpm}(x)
       =\frac{3-4x+4x^2}{3x^4(1-x)^4},\\
 G_+(x)&=G_{mmpp}(x)=G_{ppmm}(x)
       =\frac{x^2(4-4x+3x^2)}{3(1-x)^4},\\
 G_{pmmp}(x)&=G_{mppm}(x)=G_+(1-x).
 \label{expphysics:all-u-orderings}
\end{align}
The neutral sum before collecting a numerator is particularly
simple:
\[
 G_0(x)=\frac1{3x^4(1-x)^4}
       +\frac2{3x^2(1-x)^4}
       +\frac2{3x^4(1-x)^2}.
\]
The three terms correspond to the three exponent pairs in the
surviving-term table. Crossing is a rational identity:
$G_0(1-x)=G_0(x)$ and $G_+(x)=x^{-4}G_0(1/x)$.

The equal-charge OPE starts at power $x^2$. Its leading state is
the weight-six primary $\widehat u_{(-3)}\widehat u$, with the
following three momentum components:
\begin{align}
 \widehat u_{(-3)}\widehat u={}&
 -\frac{\sqrt{13}+3\sqrt{-3}}{120}
 e^{(0,0,0,1,1,-2,1,1)}\nonumber\\
 &-\frac{\sqrt{13}-3\sqrt{-3}}{120}
 e^{(1,1,0,0,0,-2,1,1)}
 -\frac1{60}e^{(1,1,0,1,1,-2,0,0)}.
 \label{expphysics:charged-leading-vector}
\end{align}
Each momentum has squared length twelve. The saved exact norm of
this state is $4/3$, agreeing with the first equal-charge global
coefficient $A_6=4/3$. This supplies an actual state behind the
leading channel, not merely a power-series coefficient.

\subsection{How to verify the global-block columns by hand}
\label{expphysics:block-recursion}

For an external unit field of weight $h$, the convention is
\begin{equation}
 R(x)=x^{2h}G(x)=A_0+\sum_{n\geq1}
 A_nx^n\,{}_2F_1(n,n;2n;x).
 \label{expphysics:global-convention}
\end{equation}
Write $R(x)=\sum_{m\geq0}r_mx^m$. Since the hypergeometric
series has leading coefficient one, this is triangular:
\begin{align}
 A_0&=r_0,\\
 A_m&=r_m-\sum_{n=1}^{m-1}A_n
 \frac{(n)_{m-n}^{\,2}}{(2n)_{m-n}(m-n)!},\qquad m\geq1.
 \label{expphysics:block-triangular}
\end{align}
Here $(n)_j$ is the rising Pochhammer symbol. This recurrence is
enough to verify every entry of the next six tables with rational
arithmetic. The $A_0$ term is constant: it has no descendant
contribution to subtract.

For example the neutral generator has $r_1=8/3$, $r_2=6$ and
$r_3=12$. The first current block contributes $A_1/2$ at order
two, so $A_2=6-4/3=14/3$. At order three it contributes $A_1/3$,
while the weight-two block contributes $A_2$, giving
$A_3=12-8/9-14/3=58/9$. A bare Taylor coefficient such as 6
therefore cannot be interpreted directly as a primary coupling sum.

In a reflection-positive channel, $A_n$ is a sum of squared
couplings to orthonormal global quasiprimaries at weight $n$.
It includes global quasiprimaries which are Virasoro descendants
of lower-weight states. It is neither the squared coupling to
one distinguished field nor the coefficient of a single
Virasoro conformal block. Odd coefficients vanish for the
identical-boson equal-charge, spin-three, current and stress
channels displayed below. The neutral charged-adjoint ordering
does not have that odd-weight selection rule.

\subsubsection{Weight-two generator: neutral channel}
The first coefficient $A_1=8/3=1/2+13/6$ resolves the two Cartan
currents, whose charges on $u$ are $(1,13)$ and whose squared
norms are $(2,78)$. At weight two the decomposition was printed
in Table~\ref{expphysics:u-pair}.
\begingroup\small
\setlength{\tabcolsep}{5pt}
\begin{longtable}{rrr}
\caption{Global-block and reduced power-series coefficients through weight 32 for the neutral weight-two generator. All entries are copied as exact fractions.}\label{expphysics:blocks-u_neutral}\\
\toprule
$n$ & $A_n$ & $[x^n]x^{2h}G(x)$\\
\midrule
\endfirsthead
\multicolumn{3}{l}{\small\itshape Continued from preceding page.}\\
\toprule
$n$ & $A_n$ & $[x^n]x^{2h}G(x)$\\
\midrule
\endhead
\midrule
\multicolumn{3}{r}{\small\itshape Continued on next page.}\\
\endfoot
\bottomrule
\endlastfoot
0 & $1$ & $1$\\
1 & $\frac{8}{3}$ & $\frac{8}{3}$\\
2 & $\frac{14}{3}$ & $6$\\
3 & $\frac{58}{9}$ & $12$\\
4 & $\frac{107}{15}$ & $\frac{65}{3}$\\
5 & $\frac{674}{105}$ & $36$\\
6 & $\frac{131}{27}$ & $56$\\
7 & $\frac{2207}{693}$ & $\frac{248}{3}$\\
8 & $\frac{2402}{1287}$ & $117$\\
9 & $\frac{19264}{19305}$ & $160$\\
10 & $\frac{18047}{36465}$ & $\frac{638}{3}$\\
11 & $\frac{31957}{138567}$ & $276$\\
12 & $\frac{53957}{529074}$ & $351$\\
13 & $\frac{12497}{289731}$ & $\frac{1316}{3}$\\
14 & $\frac{136957}{7800450}$ & $540$\\
15 & $\frac{208007}{30087450}$ & $656$\\
16 & $\frac{153811}{58169070}$ & $\frac{2363}{3}$\\
17 & $\frac{888764}{901620585}$ & $936$\\
18 & $\frac{628679}{1750204665}$ & $1102$\\
19 & $\frac{872957}{6806351475}$ & $\frac{3860}{3}$\\
20 & $\frac{170281}{3786992550}$ & $1491$\\
21 & $\frac{1603037}{103384896615}$ & $1716$\\
22 & $\frac{2126357}{403693405830}$ & $\frac{5888}{3}$\\
23 & $\frac{2785279}{1578074222790}$ & $2232$\\
24 & $\frac{901658}{1543768261425}$ & $2525$\\
25 & $\frac{4621052}{24185702762325}$ & $\frac{8528}{3}$\\
26 & $\frac{5863327}{94807954828314}$ & $3186$\\
27 & $\frac{36319}{1832211328626}$ & $3556$\\
28 & $\frac{4549}{722515868004}$ & $\frac{11861}{3}$\\
29 & $\frac{11375107}{5736517950570330}$ & $4380$\\
30 & $\frac{1996091}{3221492839236540}$ & $4836$\\
31 & $\frac{17047367}{88698436173646068}$ & $\frac{15968}{3}$\\
32 & $\frac{10333541}{174535632470722908}$ & $5841$\\
\end{longtable}
\endgroup


\subsubsection{Weight-two generator: equal-charge channel}
There is no identity or current exchange in this charge sector,
and no exchange below weight six. The leading weight-six field
is the explicit vector \eqref{expphysics:charged-leading-vector}.
\begingroup\small
\setlength{\tabcolsep}{5pt}
\begin{longtable}{rrr}
\caption{Global-block and reduced power-series coefficients through weight 32 for the equal-charge weight-two generator. All entries are copied as exact fractions.}\label{expphysics:blocks-u_equal_charge}\\
\toprule
$n$ & $A_n$ & $[x^n]x^{2h}G(x)$\\
\midrule
\endfirsthead
\multicolumn{3}{l}{\small\itshape Continued from preceding page.}\\
\toprule
$n$ & $A_n$ & $[x^n]x^{2h}G(x)$\\
\midrule
\endhead
\midrule
\multicolumn{3}{r}{\small\itshape Continued on next page.}\\
\endfoot
\bottomrule
\endlastfoot
0 & $0$ & $0$\\
1 & $0$ & $0$\\
2 & $0$ & $0$\\
3 & $0$ & $0$\\
4 & $0$ & $0$\\
5 & $0$ & $0$\\
6 & $\frac{4}{3}$ & $\frac{4}{3}$\\
7 & $0$ & $4$\\
8 & $\frac{19}{13}$ & $9$\\
9 & $0$ & $\frac{52}{3}$\\
10 & $\frac{48}{85}$ & $30$\\
11 & $0$ & $48$\\
12 & $\frac{50}{357}$ & $\frac{217}{3}$\\
13 & $0$ & $104$\\
14 & $\frac{82}{3059}$ & $144$\\
15 & $0$ & $\frac{580}{3}$\\
16 & $\frac{259}{60030}$ & $253$\\
17 & $0$ & $324$\\
18 & $\frac{6272}{10235115}$ & $\frac{1222}{3}$\\
19 & $0$ & $504$\\
20 & $\frac{724}{9147325}$ & $615$\\
21 & $0$ & $\frac{2224}{3}$\\
22 & $\frac{50}{5270837}$ & $884$\\
23 & $0$ & $1044$\\
24 & $\frac{4895}{4574128182}$ & $\frac{3667}{3}$\\
25 & $0$ & $1420$\\
26 & $\frac{6952}{60541478179}$ & $1638$\\
27 & $0$ & $\frac{5632}{3}$\\
28 & $\frac{39}{3301329773}$ & $2139$\\
29 & $0$ & $2424$\\
30 & $\frac{923}{787266089745}$ & $\frac{8200}{3}$\\
31 & $0$ & $3068$\\
32 & $\frac{85225}{755565508531268}$ & $3429$\\
\end{longtable}
\endgroup


\subsection{Spin-three self-OPEs with the residues fully specialized}

For either intrinsic spin-three algebra put $N=c/3$,
$b=a_{(1)}a$ and $\widehat a=a/\sqrt N$. The self-OPE is
\[
 a(z)a(0)\sim\frac{N}{z^6}+\frac{2T_c}{z^4}
 +\frac{\partial T_c}{z^3}+\frac b{z^2}
+\frac{\frac12\partial b-\frac1{12}\partial^3T_c}{z}.
\]
Here $a$ is the canonically rescaled generator. The raw spin-three
coordinate fields printed in Appendix~\ref{expdata:coordinates}
must first be multiplied by the conversion factors stated there.
Their raw self-OPE coefficient is not, in general, $c/3$ and is
not automatically their positive Hermitian norm. The final unit
field $\widehat a$ includes the further positive factor $1/\sqrt{c/3}$.
The two cases have
\begin{center}
\begin{tabular}{lrrrr}
\toprule
algebra&$c$&$\lambda$&$a_{(3)}b/a$&$\|b\|^2$\\
\midrule
Potts&$4/5$&$4/13$&$1119/65$&$1492/325$\\
parafermion 13&$8/5$&$14/473$&$69283/2365$&$554264/35475$\\
\bottomrule
\end{tabular}
\end{center}
These are intrinsic central charges. Replacing either with the
ambient value eight would give different stress residues.

Let $S_j(x)=x^{-j}+(1-x)^{-j}$. Specializing every coefficient,
the complete unit-normalized functions are
\begin{align}
 G_{\rm Potts}(x)={}&1+S_6(x)
 +\frac{45}{2}\{S_4(x)+S_3(x)\}
 +\frac{3357}{52}S_2(x)+\frac{1386}{13}S_1(x),\\
 G_{\rm PF13}(x)={}&1+S_6(x)
 +\frac{45}{4}\{S_4(x)+S_3(x)\}
 +\frac{207849}{3784}S_2(x)+\frac{46641}{473}S_1(x).
 \label{expphysics:spin3-specialized}
\end{align}
Their denominators can be taken to be $x^6(1-x)^6$.
For a direct polynomial comparison, write their numerators as
$P(x)=\sum_{j=0}^{12}p_jx^j$. They are palindromic,
$p_{12-j}=p_j$, and their first seven coefficients are
\begin{align*}
 (p_0,\ldots,p_6)_{\rm Potts}={}&
 \left(1,-6,\frac{75}{2},-\frac{265}{2},\frac{14667}{52},
 -\frac{5190}{13},\frac{11319}{26}\right),\\
 (p_0,\ldots,p_6)_{\rm PF13}={}&
 \left(1,-6,\frac{105}{4},-\frac{305}{4},\frac{647739}{3784},
 -\frac{138690}{473},\frac{671793}{1892}\right).
\end{align*}
The principal parts at zero follow from the displayed OPE. Equal-field
crossing fixes those at one, and the identity channel fixes the
constant at infinity. There is no remaining holomorphic function on
the sphere vanishing at infinity. This is why the rational functions
are determined beyond the finite expansion cutoff.

The intrinsic Virasoro quasiprimary
$\Lambda=:T_cT_c:-\tfrac3{10}\partial^2T_c$ has norm
$c(5c+22)/10$. Its squared unit-normalized coupling in the
$\widehat a\widehat a$ channel is $4608/[5c(5c+22)]$.
Subtracting this from the global weight-four exchange leaves
\[
 \gamma_4^2=
 \frac{48(c+2)\{8-\lambda(5c+22)\}}{c(5c+22)}
 =0\quad\hbox{or}\quad\frac{60552}{2365},
\]
for Potts and the parafermion case, respectively. The new
weight-four primary is null in Potts. Its nonzero counterpart in
the parafermion case should not be confused with the unavoidable
weight-four global quasiprimary in the Virasoro vacuum module.

\subsubsection{Potts spin-three global blocks}
\begingroup\small
\setlength{\tabcolsep}{5pt}
\begin{longtable}{rrr}
\caption{Global-block and reduced power-series coefficients through weight 32 for the Potts spin-three generator. All entries are copied as exact fractions.}\label{expphysics:blocks-Potts}\\
\toprule
$n$ & $A_n$ & $[x^n]x^{2h}G(x)$\\
\midrule
\endfirsthead
\multicolumn{3}{l}{\small\itshape Continued from preceding page.}\\
\toprule
$n$ & $A_n$ & $[x^n]x^{2h}G(x)$\\
\midrule
\endhead
\midrule
\multicolumn{3}{r}{\small\itshape Continued on next page.}\\
\endfoot
\bottomrule
\endlastfoot
0 & $1$ & $1$\\
1 & $0$ & $0$\\
2 & $\frac{45}{2}$ & $\frac{45}{2}$\\
3 & $0$ & $\frac{45}{2}$\\
4 & $\frac{576}{13}$ & $\frac{3357}{52}$\\
5 & $0$ & $\frac{1386}{13}$\\
6 & $\frac{28765}{364}$ & $\frac{11345}{52}$\\
7 & $0$ & $\frac{5190}{13}$\\
8 & $\frac{211533}{3718}$ & $\frac{35427}{52}$\\
9 & $0$ & $\frac{14246}{13}$\\
10 & $\frac{3271869}{126412}$ & $\frac{87381}{52}$\\
11 & $0$ & $\frac{32220}{13}$\\
12 & $\frac{273438}{29393}$ & $\frac{184107}{52}$\\
13 & $0$ & $\frac{64026}{13}$\\
14 & $\frac{53701083}{19315400}$ & $\frac{348381}{52}$\\
15 & $0$ & $\frac{116242}{13}$\\
16 & $\frac{7887591}{11202932}$ & $\frac{46959}{4}$\\
17 & $0$ & $\frac{197526}{13}$\\
18 & $\frac{16764239}{108734340}$ & $\frac{1009781}{52}$\\
19 & $0$ & $\frac{319032}{13}$\\
20 & $\frac{114189576}{3829070245}$ & $\frac{1596555}{52}$\\
21 & $0$ & $\frac{494826}{13}$\\
22 & $\frac{220284561}{42408196168}$ & $\frac{2433501}{52}$\\
23 & $0$ & $\frac{742302}{13}$\\
24 & $\frac{492404611}{594636663660}$ & $\frac{3597475}{52}$\\
25 & $0$ & $\frac{1082598}{13}$\\
26 & $\frac{22355809337}{182593098187864}$ & $\frac{5181141}{52}$\\
27 & $0$ & $\frac{1541012}{13}$\\
28 & $\frac{295020099}{17383387729001}$ & $\frac{7294635}{52}$\\
29 & $0$ & $165186$\\
30 & $\frac{289881678351}{130291488164677840}$ & $\frac{10067229}{52}$\\
31 & $0$ & $\frac{2936682}{13}$\\
32 & $\frac{420256422603}{1512642148079598536}$ & $\frac{13648995}{52}$\\
\end{longtable}
\endgroup

\subsubsection{Parafermion spin-three global blocks}
\begingroup\small
\setlength{\tabcolsep}{5pt}
\begin{longtable}{rrr}
\caption{Global-block and reduced power-series coefficients through weight 32 for the level-thirteen parafermion spin-three generator. All entries are copied as exact fractions.}\label{expphysics:blocks-PF13}\\
\toprule
$n$ & $A_n$ & $[x^n]x^{2h}G(x)$\\
\midrule
\endfirsthead
\multicolumn{3}{l}{\small\itshape Continued from preceding page.}\\
\toprule
$n$ & $A_n$ & $[x^n]x^{2h}G(x)$\\
\midrule
\endhead
\midrule
\multicolumn{3}{r}{\small\itshape Continued on next page.}\\
\endfoot
\bottomrule
\endlastfoot
0 & $1$ & $1$\\
1 & $0$ & $0$\\
2 & $\frac{45}{4}$ & $\frac{45}{4}$\\
3 & $0$ & $\frac{45}{4}$\\
4 & $\frac{21192}{473}$ & $\frac{207849}{3784}$\\
5 & $0$ & $\frac{46641}{473}$\\
6 & $\frac{3619235}{79464}$ & $\frac{673685}{3784}$\\
7 & $0$ & $\frac{138690}{473}$\\
8 & $\frac{9200541}{270556}$ & $\frac{1757259}{3784}$\\
9 & $0$ & $\frac{336691}{473}$\\
10 & $\frac{15229635}{836264}$ & $\frac{4017657}{3784}$\\
11 & $0$ & $\frac{731460}{473}$\\
12 & $\frac{5502941}{731731}$ & $\frac{8344119}{3784}$\\
13 & $0$ & $\frac{1459221}{473}$\\
14 & $\frac{3445217631}{1405566800}$ & $\frac{1459407}{344}$\\
15 & $0$ & $\frac{2716742}{473}$\\
16 & $\frac{1595203417}{2445686232}$ & $\frac{29007459}{3784}$\\
17 & $0$ & $\frac{4776471}{473}$\\
18 & $\frac{108115736279}{735863828040}$ & $\frac{49733777}{3784}$\\
19 & $0$ & $\frac{8001672}{473}$\\
20 & $\frac{805110996}{27863849629}$ & $\frac{81547215}{3784}$\\
21 & $0$ & $\frac{12861561}{473}$\\
22 & $\frac{172584315463}{33946129948016}$ & $\frac{128670717}{3784}$\\
23 & $0$ & $\frac{19946442}{473}$\\
24 & $\frac{1925254375}{2360250141912}$ & $\frac{196356475}{3784}$\\
25 & $0$ & $\frac{2725713}{43}$\\
26 & $\frac{1609242967069}{13287159298901488}$ & $\frac{291007017}{3784}$\\
27 & $0$ & $\frac{43848652}{473}$\\
28 & $\frac{39551299737}{2349240684519278}$ & $\frac{420296295}{3784}$\\
29 & $0$ & $\frac{62588253}{473}$\\
30 & $\frac{20967212289587}{9481211369521941280}$ & $\frac{593290773}{3784}$\\
31 & $0$ & $\frac{87427662}{473}$\\
32 & $\frac{2767264286301}{10006709594988113392}$ & $\frac{820570515}{3784}$\\
\end{longtable}
\endgroup


\subsection{Current and stress Ward correlators as normalization controls}

For a unit Cartan current and the unit total stress $\widehat T=T/2$,
\begin{align}
 G_{\widehat H}(x)&=1+x^{-2}+(1-x)^{-2},\\
 G_{\widehat T}(x)&=1+x^{-4}+(1-x)^{-4}
 +x^{-2}+(1-x)^{-2}+x^{-1}+(1-x)^{-1}.
 \label{expphysics:ward-controls}
\end{align}
These controls use the same global-block extraction and normalization
as the interacting examples. The stress exchange in the first has
$A_2=2$, whereas in the second it has $A_2=1$; their external
weights differ. The coefficient $C_{\widehat T\widehat T\widehat T}$
is one at $c=8$.

\subsubsection{Unit Cartan current}
\begingroup\small
\setlength{\tabcolsep}{5pt}
\begin{longtable}{rrr}
\caption{Global-block and reduced power-series coefficients through weight 32 for the unit Cartan current. All entries are copied as exact fractions.}\label{expphysics:blocks-unit_Cartan}\\
\toprule
$n$ & $A_n$ & $[x^n]x^{2h}G(x)$\\
\midrule
\endfirsthead
\multicolumn{3}{l}{\small\itshape Continued from preceding page.}\\
\toprule
$n$ & $A_n$ & $[x^n]x^{2h}G(x)$\\
\midrule
\endhead
\midrule
\multicolumn{3}{r}{\small\itshape Continued on next page.}\\
\endfoot
\bottomrule
\endlastfoot
0 & $1$ & $1$\\
1 & $0$ & $0$\\
2 & $2$ & $2$\\
3 & $0$ & $2$\\
4 & $\frac{6}{5}$ & $3$\\
5 & $0$ & $4$\\
6 & $\frac{5}{21}$ & $5$\\
7 & $0$ & $6$\\
8 & $\frac{14}{429}$ & $7$\\
9 & $0$ & $8$\\
10 & $\frac{9}{2431}$ & $9$\\
11 & $0$ & $10$\\
12 & $\frac{11}{29393}$ & $11$\\
13 & $0$ & $12$\\
14 & $\frac{13}{371450}$ & $13$\\
15 & $0$ & $14$\\
16 & $\frac{2}{646323}$ & $15$\\
17 & $0$ & $16$\\
18 & $\frac{17}{64822395}$ & $17$\\
19 & $0$ & $18$\\
20 & $\frac{19}{883631595}$ & $19$\\
21 & $0$ & $20$\\
22 & $\frac{7}{4077711170}$ & $21$\\
23 & $0$ & $22$\\
24 & $\frac{23}{171529806825}$ & $23$\\
25 & $0$ & $24$\\
26 & $\frac{25}{2430973200726}$ & $25$\\
27 & $0$ & $26$\\
28 & $\frac{27}{34766775458002}$ & $27$\\
29 & $0$ & $28$\\
30 & $\frac{29}{501121108325684}$ & $29$\\
31 & $0$ & $30$\\
32 & $\frac{62}{14544636039226909}$ & $31$\\
\end{longtable}
\endgroup

\subsubsection{Unit total stress}
\begingroup\small
\setlength{\tabcolsep}{5pt}
\begin{longtable}{rrr}
\caption{Global-block and reduced power-series coefficients through weight 32 for the unit total stress tensor. All entries are copied as exact fractions.}\label{expphysics:blocks-unit_total_stress}\\
\toprule
$n$ & $A_n$ & $[x^n]x^{2h}G(x)$\\
\midrule
\endfirsthead
\multicolumn{3}{l}{\small\itshape Continued from preceding page.}\\
\toprule
$n$ & $A_n$ & $[x^n]x^{2h}G(x)$\\
\midrule
\endhead
\midrule
\multicolumn{3}{r}{\small\itshape Continued on next page.}\\
\endfoot
\bottomrule
\endlastfoot
0 & $1$ & $1$\\
1 & $0$ & $0$\\
2 & $1$ & $1$\\
3 & $0$ & $1$\\
4 & $\frac{31}{10}$ & $4$\\
5 & $0$ & $7$\\
6 & $\frac{589}{126}$ & $14$\\
7 & $0$ & $25$\\
8 & $\frac{4255}{1716}$ & $41$\\
9 & $0$ & $63$\\
10 & $\frac{18569}{24310}$ & $92$\\
11 & $0$ & $129$\\
12 & $\frac{60191}{352716}$ & $175$\\
13 & $0$ & $231$\\
14 & $\frac{8439}{273700}$ & $298$\\
15 & $0$ & $377$\\
16 & $\frac{557}{116280}$ & $469$\\
17 & $0$ & $575$\\
18 & $\frac{155101}{233360622}$ & $696$\\
19 & $0$ & $833$\\
20 & $\frac{1492639}{17672631900}$ & $987$\\
21 & $0$ & $1159$\\
22 & $\frac{86851}{8681578620}$ & $1350$\\
23 & $0$ & $1561$\\
24 & $\frac{4605151}{4116715363800}$ & $1793$\\
25 & $0$ & $2047$\\
26 & $\frac{5653}{47415831372}$ & $2324$\\
27 & $0$ & $2625$\\
28 & $\frac{383105}{31402248800776}$ & $2951$\\
29 & $0$ & $3303$\\
30 & $\frac{489857}{406314412155960}$ & $3682$\\
31 & $0$ & $4089$\\
32 & $\frac{26899071}{232714176627630544}$ & $4525$\\
\end{longtable}
\endgroup


\subsection{Nine lightest-module lattice blocks and their precise scope}

The two $A_2$ lattice cosets with charges $+1$ and $-1$ restrict
to the same lightest $U$-module $M_{c_1}$, of weight $1/3$ and
top dimension three. In simple-root coordinates take
\[
 \beta_0=(2/3,1/3),\quad
 \beta_1=(-1/3,1/3),\quad
 \beta_2=(-1/3,-2/3).
\]
Their mutual Gram matrix is
\[
 ((\beta_i,\beta_j))=
 \begin{pmatrix}2/3&-1/3&-1/3\\-1/3&2/3&-1/3\\-1/3&-1/3&2/3\end{pmatrix},
\]
and their right-Cartan charges are $(2,0,-2)$. With paired lattice
intertwiners normalized to unit leading vacuum coefficient,
\begin{equation}
 \mathcal B_{ij}(x)=
 \langle V_{-\beta_i}(\infty)V_{\beta_i}(1)
          V_{-\beta_j}(x)V_{\beta_j}(0)\rangle
 =x^{-2/3}(1-x)^{-(\beta_i,\beta_j)}.
 \label{expphysics:lightest-lattice-block}
\end{equation}
Positive real powers on $0<x<1$ specify the initial branch.
Analytic continuation requires a path, since these blocks have
fractional exponents.

\begingroup\small
\setlength{\tabcolsep}{5pt}
\begin{longtable}{rrrrrr}
\caption{All nine selected lightest-module lattice blocks and their relative-weight-one decomposition. Each block also has the common factor $x^{-2/3}$.}\label{expphysics:lightest-blocks}\\
\toprule
$i$ & $j$ & $(1-x)$ exponent & total weight 1 & current part & $M_b$ part\\
\midrule
\endfirsthead
\multicolumn{6}{l}{\small\itshape Continued from preceding page.}\\
\toprule
$i$ & $j$ & $(1-x)$ exponent & total weight 1 & current part & $M_b$ part\\
\midrule
\endhead
\midrule
\multicolumn{6}{r}{\small\itshape Continued on next page.}\\
\endfoot
\bottomrule
\endlastfoot
0 & 0 & $\frac{-2}{3}$ & $\frac{2}{3}$ & $\frac{2}{39}$ & $\frac{8}{13}$\\
0 & 1 & $\frac{1}{3}$ & $\frac{-1}{3}$ & $0$ & $\frac{-1}{3}$\\
0 & 2 & $\frac{1}{3}$ & $\frac{-1}{3}$ & $\frac{-2}{39}$ & $\frac{-11}{39}$\\
1 & 0 & $\frac{1}{3}$ & $\frac{-1}{3}$ & $0$ & $\frac{-1}{3}$\\
1 & 1 & $\frac{-2}{3}$ & $\frac{2}{3}$ & $0$ & $\frac{2}{3}$\\
1 & 2 & $\frac{1}{3}$ & $\frac{-1}{3}$ & $0$ & $\frac{-1}{3}$\\
2 & 0 & $\frac{1}{3}$ & $\frac{-1}{3}$ & $\frac{-2}{39}$ & $\frac{-11}{39}$\\
2 & 1 & $\frac{1}{3}$ & $\frac{-1}{3}$ & $0$ & $\frac{-1}{3}$\\
2 & 2 & $\frac{-2}{3}$ & $\frac{2}{3}$ & $\frac{2}{39}$ & $\frac{8}{13}$\\
\end{longtable}
\endgroup


The relative-weight-one coefficient is the Gram entry
$(\beta_i,\beta_j)$. Its projection to the $U$ current is
$q_iq_j/78$. The complement in the $M_b$ top is therefore
\[
 G_b=\begin{pmatrix}
 8/13&-1/3&-11/39\\
 -1/3&2/3&-1/3\\
 -11/39&-1/3&8/13
 \end{pmatrix},\qquad
 \operatorname{spec}(G_b)=\left\{0,1,\frac{35}{39}\right\}.
\]
The zero eigenvector is $(1,1,1)$, reflecting
$\beta_0+\beta_1+\beta_2=0$. The other eigenvectors can be
chosen as $(1,-2,1)$ and $(1,0,-1)$, with eigenvalues $1$ and
$35/39$. Negative off-diagonal entries are therefore consistent
with a positive semidefinite Gram matrix.

All six permutations of $(0,1,2)$ give charge-neutral chiral
three-point functions of the three corresponding lattice
intertwiners, with absolute coefficient squared one and coordinate
factor $(z_{12}z_{13}z_{23})^{-1/3}$. Their phases depend on the
intertwiner bases and cocycle convention. The data do not provide
a complete $U$-channel basis or the sewing constants of a
diagonal-$U$ bulk four-point function. In particular, taking
$|\mathcal B_{ij}|^2$ by itself is not the claimed construction
of that nonvacuum bulk correlator.

\subsection{Exact cross-ratio samples and saved check boundaries}

These rational-point evaluations are useful when comparing an
independent implementation, plotting convention, or crossing
transformation. They are exact values, not floating-point output.

\begingroup\small
\setlength{\tabcolsep}{5pt}
\begin{longtable}{lrr}
\caption{All fifty saved rational-point evaluations. Ordering names use $p=\widehat u$, $m=\widehat u^\dagger$ at $(\infty,1,x,0)$; spin-three and Ward names refer to the unit fields.}\label{expphysics:correlator-samples}\\
\toprule
function & $x$ & exact value\\
\midrule
\endfirsthead
\multicolumn{3}{l}{\small\itshape Continued from preceding page.}\\
\toprule
function & $x$ & exact value\\
\midrule
\endhead
\midrule
\multicolumn{3}{r}{\small\itshape Continued on next page.}\\
\endfoot
\bottomrule
\endlastfoot
\texttt{ppmm} & $\frac{1}{10}$ & $\frac{121}{6561}$\\
\texttt{ppmm} & $\frac{1}{4}$ & $\frac{17}{81}$\\
\texttt{ppmm} & $\frac{1}{2}$ & $\frac{11}{3}$\\
\texttt{ppmm} & $\frac{3}{4}$ & $129$\\
\texttt{ppmm} & $\frac{9}{10}$ & $7641$\\
\texttt{pmpm} & $\frac{1}{10}$ & $\frac{88000000}{6561}$\\
\texttt{pmpm} & $\frac{1}{4}$ & $\frac{16384}{27}$\\
\texttt{pmpm} & $\frac{1}{2}$ & $\frac{512}{3}$\\
\texttt{pmpm} & $\frac{3}{4}$ & $\frac{16384}{27}$\\
\texttt{pmpm} & $\frac{9}{10}$ & $\frac{88000000}{6561}$\\
\texttt{pmmp} & $\frac{1}{10}$ & $7641$\\
\texttt{pmmp} & $\frac{1}{4}$ & $129$\\
\texttt{pmmp} & $\frac{1}{2}$ & $\frac{11}{3}$\\
\texttt{pmmp} & $\frac{3}{4}$ & $\frac{17}{81}$\\
\texttt{pmmp} & $\frac{9}{10}$ & $\frac{121}{6561}$\\
\texttt{mppm} & $\frac{1}{10}$ & $7641$\\
\texttt{mppm} & $\frac{1}{4}$ & $129$\\
\texttt{mppm} & $\frac{1}{2}$ & $\frac{11}{3}$\\
\texttt{mppm} & $\frac{3}{4}$ & $\frac{17}{81}$\\
\texttt{mppm} & $\frac{9}{10}$ & $\frac{121}{6561}$\\
\texttt{mpmp} & $\frac{1}{10}$ & $\frac{88000000}{6561}$\\
\texttt{mpmp} & $\frac{1}{4}$ & $\frac{16384}{27}$\\
\texttt{mpmp} & $\frac{1}{2}$ & $\frac{512}{3}$\\
\texttt{mpmp} & $\frac{3}{4}$ & $\frac{16384}{27}$\\
\texttt{mpmp} & $\frac{9}{10}$ & $\frac{88000000}{6561}$\\
\texttt{mmpp} & $\frac{1}{10}$ & $\frac{121}{6561}$\\
\texttt{mmpp} & $\frac{1}{4}$ & $\frac{17}{81}$\\
\texttt{mmpp} & $\frac{1}{2}$ & $\frac{11}{3}$\\
\texttt{mmpp} & $\frac{3}{4}$ & $129$\\
\texttt{mmpp} & $\frac{9}{10}$ & $7641$\\
\texttt{Potts} & $\frac{1}{10}$ & $\frac{8672450492983}{6908733}$\\
\texttt{Potts} & $\frac{1}{4}$ & $\frac{124559725}{9477}$\\
\texttt{Potts} & $\frac{1}{2}$ & $\frac{27975}{13}$\\
\texttt{Potts} & $\frac{3}{4}$ & $\frac{124559725}{9477}$\\
\texttt{Potts} & $\frac{9}{10}$ & $\frac{8672450492983}{6908733}$\\
\texttt{PF13} & $\frac{1}{10}$ & $\frac{284160941571218}{251371593}$\\
\texttt{PF13} & $\frac{1}{4}$ & $\frac{3195504725}{344817}$\\
\texttt{PF13} & $\frac{1}{2}$ & $\frac{710850}{473}$\\
\texttt{PF13} & $\frac{3}{4}$ & $\frac{3195504725}{344817}$\\
\texttt{PF13} & $\frac{9}{10}$ & $\frac{284160941571218}{251371593}$\\
\texttt{unit\_Cartan} & $\frac{1}{10}$ & $\frac{8281}{81}$\\
\texttt{unit\_Cartan} & $\frac{1}{4}$ & $\frac{169}{9}$\\
\texttt{unit\_Cartan} & $\frac{1}{2}$ & $9$\\
\texttt{unit\_Cartan} & $\frac{3}{4}$ & $\frac{169}{9}$\\
\texttt{unit\_Cartan} & $\frac{9}{10}$ & $\frac{8281}{81}$\\
\texttt{unit\_total\_stress} & $\frac{1}{10}$ & $\frac{66363661}{6561}$\\
\texttt{unit\_total\_stress} & $\frac{1}{4}$ & $\frac{22945}{81}$\\
\texttt{unit\_total\_stress} & $\frac{1}{2}$ & $45$\\
\texttt{unit\_total\_stress} & $\frac{3}{4}$ & $\frac{22945}{81}$\\
\texttt{unit\_total\_stress} & $\frac{9}{10}$ & $\frac{66363661}{6561}$\\
\end{longtable}
\endgroup


The saved report records 486 charge assignments, six $u$ orderings
with fifteen surviving terms each, exact polynomial crossing
identities for the $u$ and spin-three functions, and reconstruction
of the global expansions through weight 32. Reflection-positive
coefficients were checked through the same cutoff. Positivity at
all weights follows from the positive VOA realization when its
construction inputs hold; it is not established by extending the
finite numerical-looking table by inspection. None of these checks
is a replay of all lattice self-OPE certificates or a verification
of arbitrary full-CFT sewing identities.

\clearpage
\section{Complete modular, fusion and diagonal-defect reference tables}
\label{expphysics:defect-details}

All rows in this section use the fixed sector order
\begin{equation}
 (0,b,a,c_0,c_1,c_2,d_1,d_2,d_3,d_4,d_5,d_6).
 \label{expphysics:center-order}
\end{equation}
The identity label $0$ in a coefficient vector means the vacuum
object. A zero coefficient means its absence. Confusing these
two uses of zero is avoided by always printing a full twelve-entry
vector. These tables concern the diagonal $U$ theory; the six-line
Haagerup action after holomorphic extension is treated separately.

\subsection{Every exact modular and line eigenvalue in compact form}

For direct inspection of the cyclotomic data, put
\[
 r=\sqrt{13},\quad d=\frac{3+r}{2},\quad s=3d+2=d^2+1,
 \quad A=\frac16-\frac1{2r},\quad B=\frac16+\frac1{2r},
 \quad t_j=-2\cos\frac{2\pi j}{13}.
\]
Indices on $t_j$ are taken modulo 13 with $t_j=t_{13-j}$.
The exact $S$ matrix is completely specified by
\begin{align}
 S_{00}=S_{bb}&=A,&S_{0b}=S_{b0}&=B,\\
 S_{0x}=S_{bx}&=\frac13&& (x=a,c_0,c_1,c_2),\\
 S_{0d_j}&=r^{-1},&S_{bd_j}&=-r^{-1},\\
 S_{xd_j}&=0&& (x=a,c_0,c_1,c_2),\\
 (S_{xy})_{x,y=a,c_0,c_1,c_2}
 &=\frac13\begin{pmatrix}2&-1&-1&-1\\-1&2&-1&-1\\-1&-1&-1&2\\-1&-1&2&-1\end{pmatrix},&&\\
 S_{d_id_j}&=r^{-1}t_{ij}.&&
 \label{expphysics:exact-S}
\end{align}
Symmetry supplies the transposed entries. Thus the twelve-term
cyclotomic encodings in the JSON may be checked using a small
real trigonometric matrix. The positive embedding is essential:
$S_{00}=1/(3s)>0$.

The quantum dimensions are
\[
 (1,d^2,s,s,s,s,3d,3d,3d,3d,3d,3d),
 \qquad\mathcal D=\frac1{S_{00}}=3s.
\]
Here $\mathcal D$ is the total quantum dimension and the global
dimension is $\mathcal D^2$. For a Verlinde line $x$ the eigenvalue
on the diagonal sector $i$ is $\lambda_x(i)=S_{xi}/S_{0i}$.
For the first six sector columns, the entire eigenvalue table is
\begin{center}\small
\begin{tabular}{c|rrrrrr}
\toprule
line $x$&$0$&$b$&$a$&$c_0$&$c_1$&$c_2$\\
\midrule
$0$&1&1&1&1&1&1\\
$b$&$d^2$&$d^{-2}$&1&1&1&1\\
$a$&$s$&$s/d^2$&2&$-1$&$-1$&$-1$\\
$c_0$&$s$&$s/d^2$&$-1$&2&$-1$&$-1$\\
$c_1$&$s$&$s/d^2$&$-1$&$-1$&$-1$&2\\
$c_2$&$s$&$s/d^2$&$-1$&$-1$&2&$-1$\\
$d_i$ ($1\leq i\leq6$)&$3d$&$-3/d$&0&0&0&0\\
\bottomrule
\end{tabular}
\end{center}
For the final six sector columns, the line $0$ has six entries 1,
the line $b$ six entries $-1$, and each of $a,c_0,c_1,c_2$ six
entries 0. The remaining $6\times6$ block is
\begin{equation}
 (\lambda_{d_i}(d_j))_{i,j=1}^6=
 \begin{pmatrix}
 t_1&t_2&t_3&t_4&t_5&t_6\\
 t_2&t_4&t_6&t_5&t_3&t_1\\
 t_3&t_6&t_4&t_1&t_2&t_5\\
 t_4&t_5&t_1&t_3&t_6&t_2\\
 t_5&t_3&t_2&t_6&t_1&t_4\\
 t_6&t_1&t_5&t_2&t_4&t_3
 \end{pmatrix}.
 \label{expphysics:line-cosine-block}
\end{equation}
Equations and tables above specify all 144 exact line eigenvalues
without decimal rounding. A line eigenvalue is not a Hilbert-space
multiplicity, and may be negative or irrational.

\subsection{All center fusion products}

The full coefficient tensor is
\[
 N_{xy}^{\ \ z}=\sum_i\frac{S_{xi}S_{yi}S_{zi}}{S_{0i}}.
\]
All labels are self-dual and $S$ is real in the stated basis. The
table below prints every unordered product; commutativity gives
the remaining ordered products. For example its $b,b$ vector is
$(1,1,1,1,1,1,1,1,1,1,1,1)$, and its $a,a$ vector is
$(1,2,2,1,1,1,1,1,1,1,1,1)$. All entries are exact integers
copied from the saved Verlinde output.

\begingroup\small
\setlength{\tabcolsep}{5pt}
\begin{longtable}{rrl}
\caption{All 78 unordered center fusion products. Vectors are in the fixed order $(0,b,a,c_0,c_1,c_2,d_1,d_2,d_3,d_4,d_5,d_6)$; commutativity supplies the other 66 ordered products. Every entry of the 1728-entry tensor is therefore printed.}\label{expphysics:fusion-complete}\\
\toprule
$a$ & $b$ & coefficient vector of $a\otimes b$\\
\midrule
\endfirsthead
\multicolumn{3}{l}{\small\itshape Continued from preceding page.}\\
\toprule
$a$ & $b$ & coefficient vector of $a\otimes b$\\
\midrule
\endhead
\midrule
\multicolumn{3}{r}{\small\itshape Continued on next page.}\\
\endfoot
\bottomrule
\endlastfoot
$0$ & $0$ & $(1,0,0,0,0,0,0,0,0,0,0,0)$\\
$0$ & $b$ & $(0,1,0,0,0,0,0,0,0,0,0,0)$\\
$0$ & $a$ & $(0,0,1,0,0,0,0,0,0,0,0,0)$\\
$0$ & $c_{0}$ & $(0,0,0,1,0,0,0,0,0,0,0,0)$\\
$0$ & $c_{1}$ & $(0,0,0,0,1,0,0,0,0,0,0,0)$\\
$0$ & $c_{2}$ & $(0,0,0,0,0,1,0,0,0,0,0,0)$\\
$0$ & $d_{1}$ & $(0,0,0,0,0,0,1,0,0,0,0,0)$\\
$0$ & $d_{2}$ & $(0,0,0,0,0,0,0,1,0,0,0,0)$\\
$0$ & $d_{3}$ & $(0,0,0,0,0,0,0,0,1,0,0,0)$\\
$0$ & $d_{4}$ & $(0,0,0,0,0,0,0,0,0,1,0,0)$\\
$0$ & $d_{5}$ & $(0,0,0,0,0,0,0,0,0,0,1,0)$\\
$0$ & $d_{6}$ & $(0,0,0,0,0,0,0,0,0,0,0,1)$\\
$b$ & $b$ & $(1,1,1,1,1,1,1,1,1,1,1,1)$\\
$b$ & $a$ & $(0,1,2,1,1,1,1,1,1,1,1,1)$\\
$b$ & $c_{0}$ & $(0,1,1,2,1,1,1,1,1,1,1,1)$\\
$b$ & $c_{1}$ & $(0,1,1,1,2,1,1,1,1,1,1,1)$\\
$b$ & $c_{2}$ & $(0,1,1,1,1,2,1,1,1,1,1,1)$\\
$b$ & $d_{1}$ & $(0,1,1,1,1,1,0,1,1,1,1,1)$\\
$b$ & $d_{2}$ & $(0,1,1,1,1,1,1,0,1,1,1,1)$\\
$b$ & $d_{3}$ & $(0,1,1,1,1,1,1,1,0,1,1,1)$\\
$b$ & $d_{4}$ & $(0,1,1,1,1,1,1,1,1,0,1,1)$\\
$b$ & $d_{5}$ & $(0,1,1,1,1,1,1,1,1,1,0,1)$\\
$b$ & $d_{6}$ & $(0,1,1,1,1,1,1,1,1,1,1,0)$\\
$a$ & $a$ & $(1,2,2,1,1,1,1,1,1,1,1,1)$\\
$a$ & $c_{0}$ & $(0,1,1,1,2,2,1,1,1,1,1,1)$\\
$a$ & $c_{1}$ & $(0,1,1,2,1,2,1,1,1,1,1,1)$\\
$a$ & $c_{2}$ & $(0,1,1,2,2,1,1,1,1,1,1,1)$\\
$a$ & $d_{1}$ & $(0,1,1,1,1,1,1,1,1,1,1,1)$\\
$a$ & $d_{2}$ & $(0,1,1,1,1,1,1,1,1,1,1,1)$\\
$a$ & $d_{3}$ & $(0,1,1,1,1,1,1,1,1,1,1,1)$\\
$a$ & $d_{4}$ & $(0,1,1,1,1,1,1,1,1,1,1,1)$\\
$a$ & $d_{5}$ & $(0,1,1,1,1,1,1,1,1,1,1,1)$\\
$a$ & $d_{6}$ & $(0,1,1,1,1,1,1,1,1,1,1,1)$\\
$c_{0}$ & $c_{0}$ & $(1,2,1,2,1,1,1,1,1,1,1,1)$\\
$c_{0}$ & $c_{1}$ & $(0,1,2,1,1,2,1,1,1,1,1,1)$\\
$c_{0}$ & $c_{2}$ & $(0,1,2,1,2,1,1,1,1,1,1,1)$\\
$c_{0}$ & $d_{1}$ & $(0,1,1,1,1,1,1,1,1,1,1,1)$\\
$c_{0}$ & $d_{2}$ & $(0,1,1,1,1,1,1,1,1,1,1,1)$\\
$c_{0}$ & $d_{3}$ & $(0,1,1,1,1,1,1,1,1,1,1,1)$\\
$c_{0}$ & $d_{4}$ & $(0,1,1,1,1,1,1,1,1,1,1,1)$\\
$c_{0}$ & $d_{5}$ & $(0,1,1,1,1,1,1,1,1,1,1,1)$\\
$c_{0}$ & $d_{6}$ & $(0,1,1,1,1,1,1,1,1,1,1,1)$\\
$c_{1}$ & $c_{1}$ & $(1,2,1,1,2,1,1,1,1,1,1,1)$\\
$c_{1}$ & $c_{2}$ & $(0,1,2,2,1,1,1,1,1,1,1,1)$\\
$c_{1}$ & $d_{1}$ & $(0,1,1,1,1,1,1,1,1,1,1,1)$\\
$c_{1}$ & $d_{2}$ & $(0,1,1,1,1,1,1,1,1,1,1,1)$\\
$c_{1}$ & $d_{3}$ & $(0,1,1,1,1,1,1,1,1,1,1,1)$\\
$c_{1}$ & $d_{4}$ & $(0,1,1,1,1,1,1,1,1,1,1,1)$\\
$c_{1}$ & $d_{5}$ & $(0,1,1,1,1,1,1,1,1,1,1,1)$\\
$c_{1}$ & $d_{6}$ & $(0,1,1,1,1,1,1,1,1,1,1,1)$\\
$c_{2}$ & $c_{2}$ & $(1,2,1,1,1,2,1,1,1,1,1,1)$\\
$c_{2}$ & $d_{1}$ & $(0,1,1,1,1,1,1,1,1,1,1,1)$\\
$c_{2}$ & $d_{2}$ & $(0,1,1,1,1,1,1,1,1,1,1,1)$\\
$c_{2}$ & $d_{3}$ & $(0,1,1,1,1,1,1,1,1,1,1,1)$\\
$c_{2}$ & $d_{4}$ & $(0,1,1,1,1,1,1,1,1,1,1,1)$\\
$c_{2}$ & $d_{5}$ & $(0,1,1,1,1,1,1,1,1,1,1,1)$\\
$c_{2}$ & $d_{6}$ & $(0,1,1,1,1,1,1,1,1,1,1,1)$\\
$d_{1}$ & $d_{1}$ & $(1,0,1,1,1,1,1,0,1,1,1,1)$\\
$d_{1}$ & $d_{2}$ & $(0,1,1,1,1,1,0,1,0,1,1,1)$\\
$d_{1}$ & $d_{3}$ & $(0,1,1,1,1,1,1,0,1,0,1,1)$\\
$d_{1}$ & $d_{4}$ & $(0,1,1,1,1,1,1,1,0,1,0,1)$\\
$d_{1}$ & $d_{5}$ & $(0,1,1,1,1,1,1,1,1,0,1,0)$\\
$d_{1}$ & $d_{6}$ & $(0,1,1,1,1,1,1,1,1,1,0,0)$\\
$d_{2}$ & $d_{2}$ & $(1,0,1,1,1,1,1,1,1,0,1,1)$\\
$d_{2}$ & $d_{3}$ & $(0,1,1,1,1,1,0,1,1,1,0,1)$\\
$d_{2}$ & $d_{4}$ & $(0,1,1,1,1,1,1,0,1,1,1,0)$\\
$d_{2}$ & $d_{5}$ & $(0,1,1,1,1,1,1,1,0,1,1,0)$\\
$d_{2}$ & $d_{6}$ & $(0,1,1,1,1,1,1,1,1,0,0,1)$\\
$d_{3}$ & $d_{3}$ & $(1,0,1,1,1,1,1,1,1,1,1,0)$\\
$d_{3}$ & $d_{4}$ & $(0,1,1,1,1,1,0,1,1,1,1,0)$\\
$d_{3}$ & $d_{5}$ & $(0,1,1,1,1,1,1,0,1,1,0,1)$\\
$d_{3}$ & $d_{6}$ & $(0,1,1,1,1,1,1,1,0,0,1,1)$\\
$d_{4}$ & $d_{4}$ & $(1,0,1,1,1,1,1,1,1,1,0,1)$\\
$d_{4}$ & $d_{5}$ & $(0,1,1,1,1,1,0,1,1,0,1,1)$\\
$d_{4}$ & $d_{6}$ & $(0,1,1,1,1,1,1,0,0,1,1,1)$\\
$d_{5}$ & $d_{5}$ & $(1,0,1,1,1,1,1,1,0,1,1,1)$\\
$d_{5}$ & $d_{6}$ & $(0,1,1,1,1,1,0,0,1,1,1,1)$\\
$d_{6}$ & $d_{6}$ & $(1,0,1,1,1,1,0,1,1,1,1,1)$\\
\end{longtable}
\endgroup


One way to verify an individual row independently is to check the
dimension identity
\[
 \operatorname{qdim}(x)\operatorname{qdim}(y)
 =\sum_zN_{xy}^{\ \ z}\operatorname{qdim}(z).
\]
It is a necessary consistency check, not a replacement for the
Verlinde calculation: distinct fusion vectors can have the same
weighted dimension. The saved exporter also checks exact
associativity and all Verlinde coefficients, rather than relying
on this single scalar test.

\subsection{All twelve spatially twisted Hilbert spaces}

For a line $x$, a temporal insertion has partition function
$\sum_i\lambda_x(i)|\chi_i|^2$. Its modular transform is the
spatially twisted Hilbert-space trace
\begin{equation}
 Z_{[x]}=\sum_{i,j}(N_x)_{ij}\chi_i\overline{\chi_j},
 \qquad(N_x)_{ij}=N_{xi}^{\ \ j}.
 \label{expphysics:twisted-matrix}
\end{equation}
The matrices below have both row and column order
\eqref{expphysics:center-order}. They print all twelve complete
bilinear forms, so a reader need not transpose indices in the
fusion table while verifying a twisted spectrum. Symmetry of each
matrix follows from the real self-dual modular data.

If $h_i$ and $t_i$ denote a module's lowest weight and top dimension,
the ground dimension is the minimum of $h_i+h_j$ over positive
matrix entries. Its multiplicity is the sum of
$(N_x)_{ij}t_it_j$ over pairs attaining that minimum. Scalars
require exact equality of left and right weights, not merely
integer spin. A pair with $h_i-h_j\notin\mathbb Z$ contributes
no scalar at any grade. The vacuum module has its first positive
weight at one with six states, which must be used when excluding
the identity. These elementary rules explain why the scalar gap
can differ from the ground-state gap of a defect space.

\subsubsection*{Line $0$}
\noindent Ground dimension $0$ with multiplicity 1; first nonidentity scalar dimension $\frac{2}{3}$ with multiplicity 9.
\[\renewcommand{\arraystretch}{0.93}N_{0}=\left(\begin{array}{rrrrrrrrrrrr}
1&0&0&0&0&0&0&0&0&0&0&0\\
0&1&0&0&0&0&0&0&0&0&0&0\\
0&0&1&0&0&0&0&0&0&0&0&0\\
0&0&0&1&0&0&0&0&0&0&0&0\\
0&0&0&0&1&0&0&0&0&0&0&0\\
0&0&0&0&0&1&0&0&0&0&0&0\\
0&0&0&0&0&0&1&0&0&0&0&0\\
0&0&0&0&0&0&0&1&0&0&0&0\\
0&0&0&0&0&0&0&0&1&0&0&0\\
0&0&0&0&0&0&0&0&0&1&0&0\\
0&0&0&0&0&0&0&0&0&0&1&0\\
0&0&0&0&0&0&0&0&0&0&0&1
\end{array}\right).\]
\subsubsection*{Line $b$}
\noindent Ground dimension $\frac{2}{3}$ with multiplicity 18; first nonidentity scalar dimension $\frac{2}{3}$ with multiplicity 18.
\[\renewcommand{\arraystretch}{0.93}N_{b}=\left(\begin{array}{rrrrrrrrrrrr}
0&1&0&0&0&0&0&0&0&0&0&0\\
1&1&1&1&1&1&1&1&1&1&1&1\\
0&1&2&1&1&1&1&1&1&1&1&1\\
0&1&1&2&1&1&1&1&1&1&1&1\\
0&1&1&1&2&1&1&1&1&1&1&1\\
0&1&1&1&1&2&1&1&1&1&1&1\\
0&1&1&1&1&1&0&1&1&1&1&1\\
0&1&1&1&1&1&1&0&1&1&1&1\\
0&1&1&1&1&1&1&1&0&1&1&1\\
0&1&1&1&1&1&1&1&1&0&1&1\\
0&1&1&1&1&1&1&1&1&1&0&1\\
0&1&1&1&1&1&1&1&1&1&1&0
\end{array}\right).\]
\subsubsection*{Line $a$}
\noindent Ground dimension $\frac{2}{3}$ with multiplicity 9; first nonidentity scalar dimension $\frac{2}{3}$ with multiplicity 9.
\[\renewcommand{\arraystretch}{0.93}N_{a}=\left(\begin{array}{rrrrrrrrrrrr}
0&0&1&0&0&0&0&0&0&0&0&0\\
0&1&2&1&1&1&1&1&1&1&1&1\\
1&2&2&1&1&1&1&1&1&1&1&1\\
0&1&1&1&2&2&1&1&1&1&1&1\\
0&1&1&2&1&2&1&1&1&1&1&1\\
0&1&1&2&2&1&1&1&1&1&1&1\\
0&1&1&1&1&1&1&1&1&1&1&1\\
0&1&1&1&1&1&1&1&1&1&1&1\\
0&1&1&1&1&1&1&1&1&1&1&1\\
0&1&1&1&1&1&1&1&1&1&1&1\\
0&1&1&1&1&1&1&1&1&1&1&1\\
0&1&1&1&1&1&1&1&1&1&1&1
\end{array}\right).\]
\subsubsection*{Line $c_{0}$}
\noindent Ground dimension $\frac{2}{3}$ with multiplicity 9; first nonidentity scalar dimension $\frac{2}{3}$ with multiplicity 9.
\[\renewcommand{\arraystretch}{0.93}N_{c_{0}}=\left(\begin{array}{rrrrrrrrrrrr}
0&0&0&1&0&0&0&0&0&0&0&0\\
0&1&1&2&1&1&1&1&1&1&1&1\\
0&1&1&1&2&2&1&1&1&1&1&1\\
1&2&1&2&1&1&1&1&1&1&1&1\\
0&1&2&1&1&2&1&1&1&1&1&1\\
0&1&2&1&2&1&1&1&1&1&1&1\\
0&1&1&1&1&1&1&1&1&1&1&1\\
0&1&1&1&1&1&1&1&1&1&1&1\\
0&1&1&1&1&1&1&1&1&1&1&1\\
0&1&1&1&1&1&1&1&1&1&1&1\\
0&1&1&1&1&1&1&1&1&1&1&1\\
0&1&1&1&1&1&1&1&1&1&1&1
\end{array}\right).\]
\subsubsection*{Line $c_{1}$}
\noindent Ground dimension $\frac{1}{3}$ with multiplicity 6; first nonidentity scalar dimension $\frac{2}{3}$ with multiplicity 18.
\[\renewcommand{\arraystretch}{0.93}N_{c_{1}}=\left(\begin{array}{rrrrrrrrrrrr}
0&0&0&0&1&0&0&0&0&0&0&0\\
0&1&1&1&2&1&1&1&1&1&1&1\\
0&1&1&2&1&2&1&1&1&1&1&1\\
0&1&2&1&1&2&1&1&1&1&1&1\\
1&2&1&1&2&1&1&1&1&1&1&1\\
0&1&2&2&1&1&1&1&1&1&1&1\\
0&1&1&1&1&1&1&1&1&1&1&1\\
0&1&1&1&1&1&1&1&1&1&1&1\\
0&1&1&1&1&1&1&1&1&1&1&1\\
0&1&1&1&1&1&1&1&1&1&1&1\\
0&1&1&1&1&1&1&1&1&1&1&1\\
0&1&1&1&1&1&1&1&1&1&1&1
\end{array}\right).\]
\subsubsection*{Line $c_{2}$}
\noindent Ground dimension $\frac{2}{3}$ with multiplicity 63; first nonidentity scalar dimension $\frac{2}{3}$ with multiplicity 9.
\[\renewcommand{\arraystretch}{0.93}N_{c_{2}}=\left(\begin{array}{rrrrrrrrrrrr}
0&0&0&0&0&1&0&0&0&0&0&0\\
0&1&1&1&1&2&1&1&1&1&1&1\\
0&1&1&2&2&1&1&1&1&1&1&1\\
0&1&2&1&2&1&1&1&1&1&1&1\\
0&1&2&2&1&1&1&1&1&1&1&1\\
1&2&1&1&1&2&1&1&1&1&1&1\\
0&1&1&1&1&1&1&1&1&1&1&1\\
0&1&1&1&1&1&1&1&1&1&1&1\\
0&1&1&1&1&1&1&1&1&1&1&1\\
0&1&1&1&1&1&1&1&1&1&1&1\\
0&1&1&1&1&1&1&1&1&1&1&1\\
0&1&1&1&1&1&1&1&1&1&1&1
\end{array}\right).\]
\subsubsection*{Line $d_{1}$}
\noindent Ground dimension $\frac{6}{13}$ with multiplicity 14; first nonidentity scalar dimension $\frac{2}{3}$ with multiplicity 9.
\[\renewcommand{\arraystretch}{0.93}N_{d_{1}}=\left(\begin{array}{rrrrrrrrrrrr}
0&0&0&0&0&0&1&0&0&0&0&0\\
0&1&1&1&1&1&0&1&1&1&1&1\\
0&1&1&1&1&1&1&1&1&1&1&1\\
0&1&1&1&1&1&1&1&1&1&1&1\\
0&1&1&1&1&1&1&1&1&1&1&1\\
0&1&1&1&1&1&1&1&1&1&1&1\\
1&0&1&1&1&1&1&0&1&1&1&1\\
0&1&1&1&1&1&0&1&0&1&1&1\\
0&1&1&1&1&1&1&0&1&0&1&1\\
0&1&1&1&1&1&1&1&0&1&0&1\\
0&1&1&1&1&1&1&1&1&0&1&0\\
0&1&1&1&1&1&1&1&1&1&0&0
\end{array}\right).\]
\subsubsection*{Line $d_{2}$}
\noindent Ground dimension $\frac{2}{3}$ with multiplicity 9; first nonidentity scalar dimension $\frac{2}{3}$ with multiplicity 9.
\[\renewcommand{\arraystretch}{0.93}N_{d_{2}}=\left(\begin{array}{rrrrrrrrrrrr}
0&0&0&0&0&0&0&1&0&0&0&0\\
0&1&1&1&1&1&1&0&1&1&1&1\\
0&1&1&1&1&1&1&1&1&1&1&1\\
0&1&1&1&1&1&1&1&1&1&1&1\\
0&1&1&1&1&1&1&1&1&1&1&1\\
0&1&1&1&1&1&1&1&1&1&1&1\\
0&1&1&1&1&1&0&1&0&1&1&1\\
1&0&1&1&1&1&1&1&1&0&1&1\\
0&1&1&1&1&1&0&1&1&1&0&1\\
0&1&1&1&1&1&1&0&1&1&1&0\\
0&1&1&1&1&1&1&1&0&1&1&0\\
0&1&1&1&1&1&1&1&1&0&0&1
\end{array}\right).\]
\subsubsection*{Line $d_{3}$}
\noindent Ground dimension $\frac{2}{3}$ with multiplicity 9; first nonidentity scalar dimension $\frac{2}{3}$ with multiplicity 9.
\[\renewcommand{\arraystretch}{0.93}N_{d_{3}}=\left(\begin{array}{rrrrrrrrrrrr}
0&0&0&0&0&0&0&0&1&0&0&0\\
0&1&1&1&1&1&1&1&0&1&1&1\\
0&1&1&1&1&1&1&1&1&1&1&1\\
0&1&1&1&1&1&1&1&1&1&1&1\\
0&1&1&1&1&1&1&1&1&1&1&1\\
0&1&1&1&1&1&1&1&1&1&1&1\\
0&1&1&1&1&1&1&0&1&0&1&1\\
0&1&1&1&1&1&0&1&1&1&0&1\\
1&0&1&1&1&1&1&1&1&1&1&0\\
0&1&1&1&1&1&0&1&1&1&1&0\\
0&1&1&1&1&1&1&0&1&1&0&1\\
0&1&1&1&1&1&1&1&0&0&1&1
\end{array}\right).\]
\subsubsection*{Line $d_{4}$}
\noindent Ground dimension $\frac{5}{13}$ with multiplicity 10; first nonidentity scalar dimension $\frac{2}{3}$ with multiplicity 9.
\[\renewcommand{\arraystretch}{0.93}N_{d_{4}}=\left(\begin{array}{rrrrrrrrrrrr}
0&0&0&0&0&0&0&0&0&1&0&0\\
0&1&1&1&1&1&1&1&1&0&1&1\\
0&1&1&1&1&1&1&1&1&1&1&1\\
0&1&1&1&1&1&1&1&1&1&1&1\\
0&1&1&1&1&1&1&1&1&1&1&1\\
0&1&1&1&1&1&1&1&1&1&1&1\\
0&1&1&1&1&1&1&1&0&1&0&1\\
0&1&1&1&1&1&1&0&1&1&1&0\\
0&1&1&1&1&1&0&1&1&1&1&0\\
1&0&1&1&1&1&1&1&1&1&0&1\\
0&1&1&1&1&1&0&1&1&0&1&1\\
0&1&1&1&1&1&1&0&0&1&1&1
\end{array}\right).\]
\subsubsection*{Line $d_{5}$}
\noindent Ground dimension $\frac{7}{13}$ with multiplicity 26; first nonidentity scalar dimension $\frac{2}{3}$ with multiplicity 9.
\[\renewcommand{\arraystretch}{0.93}N_{d_{5}}=\left(\begin{array}{rrrrrrrrrrrr}
0&0&0&0&0&0&0&0&0&0&1&0\\
0&1&1&1&1&1&1&1&1&1&0&1\\
0&1&1&1&1&1&1&1&1&1&1&1\\
0&1&1&1&1&1&1&1&1&1&1&1\\
0&1&1&1&1&1&1&1&1&1&1&1\\
0&1&1&1&1&1&1&1&1&1&1&1\\
0&1&1&1&1&1&1&1&1&0&1&0\\
0&1&1&1&1&1&1&1&0&1&1&0\\
0&1&1&1&1&1&1&0&1&1&0&1\\
0&1&1&1&1&1&0&1&1&0&1&1\\
1&0&1&1&1&1&1&1&0&1&1&1\\
0&1&1&1&1&1&0&0&1&1&1&1
\end{array}\right).\]
\subsubsection*{Line $d_{6}$}
\noindent Ground dimension $\frac{8}{13}$ with multiplicity 28; first nonidentity scalar dimension $\frac{2}{3}$ with multiplicity 9.
\[\renewcommand{\arraystretch}{0.93}N_{d_{6}}=\left(\begin{array}{rrrrrrrrrrrr}
0&0&0&0&0&0&0&0&0&0&0&1\\
0&1&1&1&1&1&1&1&1&1&1&0\\
0&1&1&1&1&1&1&1&1&1&1&1\\
0&1&1&1&1&1&1&1&1&1&1&1\\
0&1&1&1&1&1&1&1&1&1&1&1\\
0&1&1&1&1&1&1&1&1&1&1&1\\
0&1&1&1&1&1&1&1&1&1&0&0\\
0&1&1&1&1&1&1&1&1&0&0&1\\
0&1&1&1&1&1&1&1&0&0&1&1\\
0&1&1&1&1&1&1&0&0&1&1&1\\
0&1&1&1&1&1&0&0&1&1&1&1\\
1&0&1&1&1&1&0&1&1&1&1&1
\end{array}\right).\]


\subsection{All diagonal Cardy annuli and their positive gaps}

The elementary boundary labeled $x$ has closed-channel state
\[
 |x\rangle=\sum_i\frac{S_{xi}}{\sqrt{S_{0i}}}|i\rangle\!\rangle.
\]
The open annulus between $x$ and $y$ is
$Z_{xy}^{\rm open}=\sum_iN_{xy}^{\ \ i}\chi_i$. Accordingly
Table~\ref{expphysics:fusion-complete} already prints every full
annulus character vector. The next table adds all ground-state
and first-positive-weight data. A self-annulus has one vacuum,
whereas distinct elementary boundaries have none. The first
positive weight in a vacuum character is one, not an absent
state obtained by deleting the entire vacuum representation.

\begingroup\small
\setlength{\tabcolsep}{5pt}
\begin{longtable}{rrrrrr}
\caption{Ground states and first positive boundary weights for all 78 distinct Cardy annuli. The complete character coefficient vector is the corresponding row of Table~\ref{expphysics:fusion-complete}. Reversing the boundaries gives the same row.}\label{expphysics:cardy-gaps}\\
\toprule
$a$ & $b$ & $h_{\min}$ & vacua & $h_{>0}$ & multiplicity\\
\midrule
\endfirsthead
\multicolumn{6}{l}{\small\itshape Continued from preceding page.}\\
\toprule
$a$ & $b$ & $h_{\min}$ & vacua & $h_{>0}$ & multiplicity\\
\midrule
\endhead
\midrule
\multicolumn{6}{r}{\small\itshape Continued on next page.}\\
\endfoot
\bottomrule
\endlastfoot
$0$ & $0$ & $0$ & 1 & $1$ & 6\\
$0$ & $b$ & $1$ & 0 & $1$ & 80\\
$0$ & $a$ & $1$ & 0 & $1$ & 81\\
$0$ & $c_{0}$ & $1$ & 0 & $1$ & 81\\
$0$ & $c_{1}$ & $\frac{1}{3}$ & 0 & $\frac{1}{3}$ & 3\\
$0$ & $c_{2}$ & $\frac{2}{3}$ & 0 & $\frac{2}{3}$ & 27\\
$0$ & $d_{1}$ & $\frac{6}{13}$ & 0 & $\frac{6}{13}$ & 7\\
$0$ & $d_{2}$ & $\frac{11}{13}$ & 0 & $\frac{11}{13}$ & 42\\
$0$ & $d_{3}$ & $\frac{15}{13}$ & 0 & $\frac{15}{13}$ & 119\\
$0$ & $d_{4}$ & $\frac{5}{13}$ & 0 & $\frac{5}{13}$ & 5\\
$0$ & $d_{5}$ & $\frac{7}{13}$ & 0 & $\frac{7}{13}$ & 13\\
$0$ & $d_{6}$ & $\frac{8}{13}$ & 0 & $\frac{8}{13}$ & 14\\
$b$ & $b$ & $0$ & 1 & $\frac{1}{3}$ & 3\\
$b$ & $a$ & $\frac{1}{3}$ & 0 & $\frac{1}{3}$ & 3\\
$b$ & $c_{0}$ & $\frac{1}{3}$ & 0 & $\frac{1}{3}$ & 3\\
$b$ & $c_{1}$ & $\frac{1}{3}$ & 0 & $\frac{1}{3}$ & 6\\
$b$ & $c_{2}$ & $\frac{1}{3}$ & 0 & $\frac{1}{3}$ & 3\\
$b$ & $d_{1}$ & $\frac{1}{3}$ & 0 & $\frac{1}{3}$ & 3\\
$b$ & $d_{2}$ & $\frac{1}{3}$ & 0 & $\frac{1}{3}$ & 3\\
$b$ & $d_{3}$ & $\frac{1}{3}$ & 0 & $\frac{1}{3}$ & 3\\
$b$ & $d_{4}$ & $\frac{1}{3}$ & 0 & $\frac{1}{3}$ & 3\\
$b$ & $d_{5}$ & $\frac{1}{3}$ & 0 & $\frac{1}{3}$ & 3\\
$b$ & $d_{6}$ & $\frac{1}{3}$ & 0 & $\frac{1}{3}$ & 3\\
$a$ & $a$ & $0$ & 1 & $\frac{1}{3}$ & 3\\
$a$ & $c_{0}$ & $\frac{1}{3}$ & 0 & $\frac{1}{3}$ & 6\\
$a$ & $c_{1}$ & $\frac{1}{3}$ & 0 & $\frac{1}{3}$ & 3\\
$a$ & $c_{2}$ & $\frac{1}{3}$ & 0 & $\frac{1}{3}$ & 6\\
$a$ & $d_{1}$ & $\frac{1}{3}$ & 0 & $\frac{1}{3}$ & 3\\
$a$ & $d_{2}$ & $\frac{1}{3}$ & 0 & $\frac{1}{3}$ & 3\\
$a$ & $d_{3}$ & $\frac{1}{3}$ & 0 & $\frac{1}{3}$ & 3\\
$a$ & $d_{4}$ & $\frac{1}{3}$ & 0 & $\frac{1}{3}$ & 3\\
$a$ & $d_{5}$ & $\frac{1}{3}$ & 0 & $\frac{1}{3}$ & 3\\
$a$ & $d_{6}$ & $\frac{1}{3}$ & 0 & $\frac{1}{3}$ & 3\\
$c_{0}$ & $c_{0}$ & $0$ & 1 & $\frac{1}{3}$ & 3\\
$c_{0}$ & $c_{1}$ & $\frac{1}{3}$ & 0 & $\frac{1}{3}$ & 3\\
$c_{0}$ & $c_{2}$ & $\frac{1}{3}$ & 0 & $\frac{1}{3}$ & 6\\
$c_{0}$ & $d_{1}$ & $\frac{1}{3}$ & 0 & $\frac{1}{3}$ & 3\\
$c_{0}$ & $d_{2}$ & $\frac{1}{3}$ & 0 & $\frac{1}{3}$ & 3\\
$c_{0}$ & $d_{3}$ & $\frac{1}{3}$ & 0 & $\frac{1}{3}$ & 3\\
$c_{0}$ & $d_{4}$ & $\frac{1}{3}$ & 0 & $\frac{1}{3}$ & 3\\
$c_{0}$ & $d_{5}$ & $\frac{1}{3}$ & 0 & $\frac{1}{3}$ & 3\\
$c_{0}$ & $d_{6}$ & $\frac{1}{3}$ & 0 & $\frac{1}{3}$ & 3\\
$c_{1}$ & $c_{1}$ & $0$ & 1 & $\frac{1}{3}$ & 6\\
$c_{1}$ & $c_{2}$ & $\frac{1}{3}$ & 0 & $\frac{1}{3}$ & 3\\
$c_{1}$ & $d_{1}$ & $\frac{1}{3}$ & 0 & $\frac{1}{3}$ & 3\\
$c_{1}$ & $d_{2}$ & $\frac{1}{3}$ & 0 & $\frac{1}{3}$ & 3\\
$c_{1}$ & $d_{3}$ & $\frac{1}{3}$ & 0 & $\frac{1}{3}$ & 3\\
$c_{1}$ & $d_{4}$ & $\frac{1}{3}$ & 0 & $\frac{1}{3}$ & 3\\
$c_{1}$ & $d_{5}$ & $\frac{1}{3}$ & 0 & $\frac{1}{3}$ & 3\\
$c_{1}$ & $d_{6}$ & $\frac{1}{3}$ & 0 & $\frac{1}{3}$ & 3\\
$c_{2}$ & $c_{2}$ & $0$ & 1 & $\frac{1}{3}$ & 3\\
$c_{2}$ & $d_{1}$ & $\frac{1}{3}$ & 0 & $\frac{1}{3}$ & 3\\
$c_{2}$ & $d_{2}$ & $\frac{1}{3}$ & 0 & $\frac{1}{3}$ & 3\\
$c_{2}$ & $d_{3}$ & $\frac{1}{3}$ & 0 & $\frac{1}{3}$ & 3\\
$c_{2}$ & $d_{4}$ & $\frac{1}{3}$ & 0 & $\frac{1}{3}$ & 3\\
$c_{2}$ & $d_{5}$ & $\frac{1}{3}$ & 0 & $\frac{1}{3}$ & 3\\
$c_{2}$ & $d_{6}$ & $\frac{1}{3}$ & 0 & $\frac{1}{3}$ & 3\\
$d_{1}$ & $d_{1}$ & $0$ & 1 & $\frac{1}{3}$ & 3\\
$d_{1}$ & $d_{2}$ & $\frac{1}{3}$ & 0 & $\frac{1}{3}$ & 3\\
$d_{1}$ & $d_{3}$ & $\frac{1}{3}$ & 0 & $\frac{1}{3}$ & 3\\
$d_{1}$ & $d_{4}$ & $\frac{1}{3}$ & 0 & $\frac{1}{3}$ & 3\\
$d_{1}$ & $d_{5}$ & $\frac{1}{3}$ & 0 & $\frac{1}{3}$ & 3\\
$d_{1}$ & $d_{6}$ & $\frac{1}{3}$ & 0 & $\frac{1}{3}$ & 3\\
$d_{2}$ & $d_{2}$ & $0$ & 1 & $\frac{1}{3}$ & 3\\
$d_{2}$ & $d_{3}$ & $\frac{1}{3}$ & 0 & $\frac{1}{3}$ & 3\\
$d_{2}$ & $d_{4}$ & $\frac{1}{3}$ & 0 & $\frac{1}{3}$ & 3\\
$d_{2}$ & $d_{5}$ & $\frac{1}{3}$ & 0 & $\frac{1}{3}$ & 3\\
$d_{2}$ & $d_{6}$ & $\frac{1}{3}$ & 0 & $\frac{1}{3}$ & 3\\
$d_{3}$ & $d_{3}$ & $0$ & 1 & $\frac{1}{3}$ & 3\\
$d_{3}$ & $d_{4}$ & $\frac{1}{3}$ & 0 & $\frac{1}{3}$ & 3\\
$d_{3}$ & $d_{5}$ & $\frac{1}{3}$ & 0 & $\frac{1}{3}$ & 3\\
$d_{3}$ & $d_{6}$ & $\frac{1}{3}$ & 0 & $\frac{1}{3}$ & 3\\
$d_{4}$ & $d_{4}$ & $0$ & 1 & $\frac{1}{3}$ & 3\\
$d_{4}$ & $d_{5}$ & $\frac{1}{3}$ & 0 & $\frac{1}{3}$ & 3\\
$d_{4}$ & $d_{6}$ & $\frac{1}{3}$ & 0 & $\frac{1}{3}$ & 3\\
$d_{5}$ & $d_{5}$ & $0$ & 1 & $\frac{1}{3}$ & 3\\
$d_{5}$ & $d_{6}$ & $\frac{1}{3}$ & 0 & $\frac{1}{3}$ & 3\\
$d_{6}$ & $d_{6}$ & $0$ & 1 & $\frac{1}{3}$ & 3\\
\end{longtable}
\endgroup


Boundary overlaps in this normalization are
\begin{equation}
 g_x=\frac{\operatorname{qdim}(x)}{\sqrt{3s}},\qquad
 g_0^2=\frac1{3s},\quad g_b^2=\frac{d^4}{3s},\quad
 g_a^2=g_{c_j}^2=\frac{s}{3},\quad
 g_{d_j}^2=\frac{3d^2}{s}.
 \label{expphysics:cardy-g-exact}
\end{equation}
Their positive values, shown with the saved plotting precision, are
\begin{center}
\begin{tabular}{lr}
\toprule
boundary type&$g_x$\\\midrule
$0$&$0.1673069562521514$\\
$b$&$1.825038973650249$\\
$a,c_j$&$1.992345929902401$\\
$d_j$&$1.657732017398097$\\
\bottomrule
\end{tabular}
\end{center}
These decimals are the saved plotting values; the displayed
algebraic expressions define the observables. The boundary entropy
is $\log g_x$, not $g_x$ itself.

The vacuum Cardy boundary has six positive-weight-one currents.
Every other self-annulus contains weight-$1/3$ boundary fields:
six for $c_1$ and three for each other nonvacuum boundary. This
identifies relevant boundary operators. Neither their degeneracy
nor their conformal weight specifies a boundary renormalization
group flow or its endpoint.

\clearpage
\section{The six Haagerup lines and all extension boundary sectors}
\label{expphysics:extension-details}

The selected holomorphic extension is $H'=V_{E_8}$ with
$H'|_U=M_0\oplus M_b\oplus2M_a$. Its six right-module labels are
ordered as
\[
 (1,\alpha,\alpha^2,\rho,\alpha\rho,\alpha^2\rho).
\]
They have dimensions $(1,1,1,d,d,d)$ and fusion rules
$\alpha^3=1$, $\rho\alpha=\alpha^2\rho$,
$\rho^2=1+\rho+\alpha\rho+\alpha^2\rho$. Unlike the center
fusion table, the ordered product table need not be symmetric.

\begingroup\small
\setlength{\tabcolsep}{5pt}
\begin{longtable}{rrl}
\caption{All 36 ordered Haagerup products, in basis $(1,\alpha,\alpha^2,\rho,\alpha\rho,\alpha^2\rho)$. Order matters.}\label{expphysics:h3-fusion}\\
\toprule
$x$ & $y$ & coefficient vector of $x\otimes y$\\
\midrule
\endfirsthead
\multicolumn{3}{l}{\small\itshape Continued from preceding page.}\\
\toprule
$x$ & $y$ & coefficient vector of $x\otimes y$\\
\midrule
\endhead
\midrule
\multicolumn{3}{r}{\small\itshape Continued on next page.}\\
\endfoot
\bottomrule
\endlastfoot
$1$ & $1$ & $(1,0,0,0,0,0)$\\
$1$ & $\alpha$ & $(0,1,0,0,0,0)$\\
$1$ & $\alpha^2$ & $(0,0,1,0,0,0)$\\
$1$ & $\rho$ & $(0,0,0,1,0,0)$\\
$1$ & $\alpha\rho$ & $(0,0,0,0,1,0)$\\
$1$ & $\alpha^2\rho$ & $(0,0,0,0,0,1)$\\
$\alpha$ & $1$ & $(0,1,0,0,0,0)$\\
$\alpha$ & $\alpha$ & $(0,0,1,0,0,0)$\\
$\alpha$ & $\alpha^2$ & $(1,0,0,0,0,0)$\\
$\alpha$ & $\rho$ & $(0,0,0,0,1,0)$\\
$\alpha$ & $\alpha\rho$ & $(0,0,0,0,0,1)$\\
$\alpha$ & $\alpha^2\rho$ & $(0,0,0,1,0,0)$\\
$\alpha^2$ & $1$ & $(0,0,1,0,0,0)$\\
$\alpha^2$ & $\alpha$ & $(1,0,0,0,0,0)$\\
$\alpha^2$ & $\alpha^2$ & $(0,1,0,0,0,0)$\\
$\alpha^2$ & $\rho$ & $(0,0,0,0,0,1)$\\
$\alpha^2$ & $\alpha\rho$ & $(0,0,0,1,0,0)$\\
$\alpha^2$ & $\alpha^2\rho$ & $(0,0,0,0,1,0)$\\
$\rho$ & $1$ & $(0,0,0,1,0,0)$\\
$\rho$ & $\alpha$ & $(0,0,0,0,0,1)$\\
$\rho$ & $\alpha^2$ & $(0,0,0,0,1,0)$\\
$\rho$ & $\rho$ & $(1,0,0,1,1,1)$\\
$\rho$ & $\alpha\rho$ & $(0,0,1,1,1,1)$\\
$\rho$ & $\alpha^2\rho$ & $(0,1,0,1,1,1)$\\
$\alpha\rho$ & $1$ & $(0,0,0,0,1,0)$\\
$\alpha\rho$ & $\alpha$ & $(0,0,0,1,0,0)$\\
$\alpha\rho$ & $\alpha^2$ & $(0,0,0,0,0,1)$\\
$\alpha\rho$ & $\rho$ & $(0,1,0,1,1,1)$\\
$\alpha\rho$ & $\alpha\rho$ & $(1,0,0,1,1,1)$\\
$\alpha\rho$ & $\alpha^2\rho$ & $(0,0,1,1,1,1)$\\
$\alpha^2\rho$ & $1$ & $(0,0,0,0,0,1)$\\
$\alpha^2\rho$ & $\alpha$ & $(0,0,0,0,1,0)$\\
$\alpha^2\rho$ & $\alpha^2$ & $(0,0,0,1,0,0)$\\
$\alpha^2\rho$ & $\rho$ & $(0,0,1,1,1,1)$\\
$\alpha^2\rho$ & $\alpha\rho$ & $(0,1,0,1,1,1)$\\
$\alpha^2\rho$ & $\alpha^2\rho$ & $(1,0,0,1,1,1)$\\
\end{longtable}
\endgroup


\subsection{Forgetful matrix, tube factors and first-factor traces}

The entire forgetful multiplicity matrix, with the six rows above
and twelve columns in \eqref{expphysics:center-order}, is
\begin{equation}
 n=\left(\begin{array}{rrrrrrrrrrrr}
 1&1&2&0&0&0&0&0&0&0&0&0\\
 0&0&0&1&1&1&0&0&0&0&0&0\\
 0&0&0&1&1&1&0&0&0&0&0&0\\
 0&1&1&1&1&1&1&1&1&1&1&1\\
 0&1&1&1&1&1&1&1&1&1&1&1\\
 0&1&1&1&1&1&1&1&1&1&1&1
 \end{array}\right).
 \label{expphysics:forgetful-matrix}
\end{equation}
Restriction of a right module $x$ is
$Rx=\bigoplus_\mu n_{x\mu}M_\mu$. The selected action is on
the first chiral factor, with the antiholomorphic $H'$ theory a
spectator. Its Hilbert space and tube decomposition are
\begin{align}
 \mathcal H_x&=Rx\otimes\overline{H'}
 =\bigoplus_\mu W_x^\mu\otimes\mathcal V_\mu,\\
 \dim W_x^\mu&=n_{x\mu},\qquad
 \mathcal V_\mu=M_\mu\otimes\overline{H'},\\
 Z_x&=\sum_\mu n_{x\mu}\chi_\mu\overline{\chi_{E_8}}.
 \label{expphysics:tube-conventions}
\end{align}
The twelve functions $\chi_\mu\overline{\chi_{E_8}}$ count the
$\mathcal V_\mu$ factors; they exclude $W_x^\mu$.
For example $\dim W_1^a=2$, and omitting that factor would
incorrectly halve the corresponding physical untwisted charged
multiplicity. Conversely inserting it into the bootstrap
multiplicity-space function would count that factor twice.

\begingroup\small
\setlength{\tabcolsep}{5pt}
\begin{longtable}{rrrrr}
\caption{All twelve tube multiplicity spaces $\mathcal V_\mu=M_\mu\otimes\overline{V_{E_8}}$. These rows exclude the finite-dimensional $W_x^\mu$ factor. ``None means there is no scalar at any grade.}\label{expphysics:tube-spaces}\\
\toprule
$\mu$ & $\Delta_{\min}$ & mult. & $\Delta_{\rm scalar}$ & mult.\\
\midrule
\endfirsthead
\multicolumn{5}{l}{\small\itshape Continued from preceding page.}\\
\toprule
$\mu$ & $\Delta_{\min}$ & mult. & $\Delta_{\rm scalar}$ & mult.\\
\midrule
\endhead
\midrule
\multicolumn{5}{r}{\small\itshape Continued on next page.}\\
\endfoot
\bottomrule
\endlastfoot
$0$ & $0$ & 1 & $2$ & 1488\\
$b$ & $1$ & 80 & $2$ & 19840\\
$a$ & $1$ & 81 & $2$ & 20088\\
$c_{0}$ & $1$ & 81 & $2$ & 20088\\
$c_{1}$ & $\frac{1}{3}$ & 3 & none & 0\\
$c_{2}$ & $\frac{2}{3}$ & 27 & none & 0\\
$d_{1}$ & $\frac{6}{13}$ & 7 & none & 0\\
$d_{2}$ & $\frac{11}{13}$ & 42 & none & 0\\
$d_{3}$ & $\frac{15}{13}$ & 119 & none & 0\\
$d_{4}$ & $\frac{5}{13}$ & 5 & none & 0\\
$d_{5}$ & $\frac{7}{13}$ & 13 & none & 0\\
$d_{6}$ & $\frac{8}{13}$ & 14 & none & 0\\
\end{longtable}
\endgroup


All antiholomorphic weights of $H'$ are integers. A sector with
noninteger $h_\mu$ therefore has no scalar at any grade, because
adding integral grades never makes its left weight equal an integer
right weight. In the four integer-weight sectors, the first
nonidentity scalars occur at $(1,1)$ and have multiplicities
$6\cdot248$, $80\cdot248$, $81\cdot248$, $81\cdot248$.
The resulting physical six twisted spaces, including the
forgetful multiplicities, are the following.

\begingroup\small
\setlength{\tabcolsep}{5pt}
\begin{longtable}{rrrrr}
\caption{All six first-factor Haagerup twisted Hilbert spaces, with the full antiholomorphic $E_8$ theory as spectator. Multiplicities here include the forgetful matrix.}\label{expphysics:h3-twisted-gaps}\\
\toprule
$x$ & $\Delta_{\min}$ & mult. & $\Delta_{\rm scalar}$ & mult.\\
\midrule
\endfirsthead
\multicolumn{5}{l}{\small\itshape Continued from preceding page.}\\
\toprule
$x$ & $\Delta_{\min}$ & mult. & $\Delta_{\rm scalar}$ & mult.\\
\midrule
\endhead
\midrule
\multicolumn{5}{r}{\small\itshape Continued on next page.}\\
\endfoot
\bottomrule
\endlastfoot
$1$ & $0$ & 1 & $2$ & 61504\\
$\alpha$ & $\frac{1}{3}$ & 3 & $2$ & 20088\\
$\alpha^2$ & $\frac{1}{3}$ & 3 & $2$ & 20088\\
$\rho$ & $\frac{1}{3}$ & 3 & $2$ & 60016\\
$\alpha\rho$ & $\frac{1}{3}$ & 3 & $2$ & 60016\\
$\alpha^2\rho$ & $\frac{1}{3}$ & 3 & $2$ & 60016\\
\end{longtable}
\endgroup


For orientation, the untwisted invariant scalar multiplicity is
1488, while the entire untwisted scalar space at dimension two
has multiplicity $248^2=61504$. The untwisted $\mathbb Z_3$-charged
space contains two copies of $\mathcal V_a$, so its scalar
multiplicity is $2\cdot81\cdot248=40176$. In contrast, the
bootstrap union $\mathcal V_a\cup\mathcal V_{c_0}$ concerns
two tube multiplicity-space labels and is not that entire charged
subspace. Their identical scalar gap two does not identify the
underlying spaces.

On the three untwisted multiplicity factors, the line matrices
may be chosen as
\[
 \begin{array}{c|cc}
  &\alpha&\rho\\\hline
 M_0&1&d\\
 M_b&1&-d^{-1}\\
 \mathbb C^2\otimes M_a&
 \begin{pmatrix}\omega&0\\0&\omega^2\end{pmatrix}&
 \begin{pmatrix}0&1\\1&0\end{pmatrix}
 \end{array},\qquad\omega=e^{2\pi i/3}.
\]
Taking the traces on these multiplicity factors gives
\[
 Z^{[\alpha]}=(\chi_0+\chi_b-\chi_a)\overline{\chi_{E_8}},
 \qquad
 Z^{[\rho]}=(d\chi_0-d^{-1}\chi_b)\overline{\chi_{E_8}}.
\]
The antiholomorphic character is unchanged in a first-factor
insertion. Replacing either expression with the absolute square
of its holomorphic factor would describe a different paired
left/right insertion.

\subsection{All \texorpdfstring{$U$}{U}-preserving annuli of the full extension}

The open Hilbert space between boundary modules $x,y$ is
$R(x^*\otimes_Ay)$. Thus
\begin{equation}
 Z_{xy}^{E_8,\mathrm{open}}=
 \sum_{z,\mu}N_{x^*y}^{\ \ z}n_{z\mu}\chi_\mu.
 \label{expphysics:extension-annulus}
\end{equation}
The duals are $1^*=1$, $\alpha^*=\alpha^2$,
$(\alpha^2)^*=\alpha$ and the three noninvertible objects
are self-dual. The dual on $x$ cannot be dropped just because
the center labels in the earlier Cardy table were self-dual.
For instance the $\alpha,\alpha$ annulus has a vacuum;
using $\alpha\otimes\alpha$ instead of
$\alpha^*\otimes\alpha$ would incorrectly remove it.

\begingroup\footnotesize
\setlength{\tabcolsep}{5pt}
\begin{longtable}{rrlrr}
\caption{All 21 distinct $U$-preserving annuli of the selected full $E_8$ theory. The vector order is the twelve-sector center order. Equal boundaries have one identity; distinct boundaries have none. Reversal gives the other 15 ordered annuli.}\label{expphysics:e8-annuli}\\
\toprule
$x$ & $y$ & $U$-character coefficient vector & $h_{>0}$ & mult.\\
\midrule
\endfirsthead
\multicolumn{5}{l}{\small\itshape Continued from preceding page.}\\
\toprule
$x$ & $y$ & $U$-character coefficient vector & $h_{>0}$ & mult.\\
\midrule
\endhead
\midrule
\multicolumn{5}{r}{\small\itshape Continued on next page.}\\
\endfoot
\bottomrule
\endlastfoot
$1$ & $1$ & $(1,1,2,0,0,0,0,0,0,0,0,0)$ & $1$ & 248\\
$1$ & $\alpha$ & $(0,0,0,1,1,1,0,0,0,0,0,0)$ & $\frac{1}{3}$ & 3\\
$1$ & $\alpha^2$ & $(0,0,0,1,1,1,0,0,0,0,0,0)$ & $\frac{1}{3}$ & 3\\
$1$ & $\rho$ & $(0,1,1,1,1,1,1,1,1,1,1,1)$ & $\frac{1}{3}$ & 3\\
$1$ & $\alpha\rho$ & $(0,1,1,1,1,1,1,1,1,1,1,1)$ & $\frac{1}{3}$ & 3\\
$1$ & $\alpha^2\rho$ & $(0,1,1,1,1,1,1,1,1,1,1,1)$ & $\frac{1}{3}$ & 3\\
$\alpha$ & $\alpha$ & $(1,1,2,0,0,0,0,0,0,0,0,0)$ & $1$ & 248\\
$\alpha$ & $\alpha^2$ & $(0,0,0,1,1,1,0,0,0,0,0,0)$ & $\frac{1}{3}$ & 3\\
$\alpha$ & $\rho$ & $(0,1,1,1,1,1,1,1,1,1,1,1)$ & $\frac{1}{3}$ & 3\\
$\alpha$ & $\alpha\rho$ & $(0,1,1,1,1,1,1,1,1,1,1,1)$ & $\frac{1}{3}$ & 3\\
$\alpha$ & $\alpha^2\rho$ & $(0,1,1,1,1,1,1,1,1,1,1,1)$ & $\frac{1}{3}$ & 3\\
$\alpha^2$ & $\alpha^2$ & $(1,1,2,0,0,0,0,0,0,0,0,0)$ & $1$ & 248\\
$\alpha^2$ & $\rho$ & $(0,1,1,1,1,1,1,1,1,1,1,1)$ & $\frac{1}{3}$ & 3\\
$\alpha^2$ & $\alpha\rho$ & $(0,1,1,1,1,1,1,1,1,1,1,1)$ & $\frac{1}{3}$ & 3\\
$\alpha^2$ & $\alpha^2\rho$ & $(0,1,1,1,1,1,1,1,1,1,1,1)$ & $\frac{1}{3}$ & 3\\
$\rho$ & $\rho$ & $(1,4,5,3,3,3,3,3,3,3,3,3)$ & $\frac{1}{3}$ & 9\\
$\rho$ & $\alpha\rho$ & $(0,3,3,4,4,4,3,3,3,3,3,3)$ & $\frac{1}{3}$ & 12\\
$\rho$ & $\alpha^2\rho$ & $(0,3,3,4,4,4,3,3,3,3,3,3)$ & $\frac{1}{3}$ & 12\\
$\alpha\rho$ & $\alpha\rho$ & $(1,4,5,3,3,3,3,3,3,3,3,3)$ & $\frac{1}{3}$ & 9\\
$\alpha\rho$ & $\alpha^2\rho$ & $(0,3,3,4,4,4,3,3,3,3,3,3)$ & $\frac{1}{3}$ & 12\\
$\alpha^2\rho$ & $\alpha^2\rho$ & $(1,4,5,3,3,3,3,3,3,3,3,3)$ & $\frac{1}{3}$ & 9\\
\end{longtable}
\endgroup


The boundary overlaps are exactly
\[
 (g_1,g_\alpha,g_{\alpha^2},g_\rho,g_{\alpha\rho},g_{\alpha^2\rho})
 =(1,1,1,d,d,d).
\]
This normalization is fixed by the vacuum coefficient in the
closed channel, rather than by an arbitrary proportionality
constant. Restriction multiplies categorical dimensions by
$\operatorname{qdim}(A)=\mathcal D$, so multiplying by
$S_{00}=1/\mathcal D$ gives
$\operatorname{qdim}(x)\operatorname{qdim}(y)$ for that coefficient.
Positive boundary overlaps then give $g_x=\operatorname{qdim}(x)$.

The three invertible self-annuli equal $\chi_{E_8}$ and have
248 weight-one currents as their first positive boundary states.
Each noninvertible self-annulus instead has nine weight-$1/3$
states. Distinct noninvertible boundaries have twelve
weight-$1/3$ ground fields between them. All other distinct pairs
have three such ground fields. These are consequences of the
printed vectors and the $M_{c_1}$ top dimension three.

\subsection{Provenance and what has actually been incorporated}

The modular, fusion, defect, and boundary entries above are saved
exact consequences of the positive Haagerup modular data and the
supplied $\gamma=0$ character tops. The source exporter represents
numbers in $\mathbb Q(\zeta_{13})$ by twelve rational coefficients
and uses the positive Gauss sum
\[
 \sqrt{13}=\sum_{j=1}^{12}\left(\frac j{13}\right)\zeta_{13}^j.
\]
It records checks of $S^2=1$, $(S\Theta)^3=1$, nonnegative
integrality of the complete Verlinde tensor, associativity,
dimension multiplicativity, self-duality, and the forgetful
matrix's compatibility with the two fusion rings. These checks
do not reconstruct associators, braidings, or generic nonvacuum
OPE coefficients from modular data.

For reproducibility, the new transcription utility is
\artifact{physics_tables/build_physics_tables.py}; its input
hashes and coverage counts are recorded in
\artifact{physics_tables/transcription_manifest.json} and
\artifact{physics_sources.md}. It reads saved finite outputs and
formats exact entries. It does not import the lattice engine or
rerun a physics exporter. Its coverage checks include all 113
component labels, the 318 distinct nonzero multiplet triples,
all 18,365 sparse component entries when identifying that support,
198 global-block rows, all center fusion products, and every
boundary pair up to the stated reversal symmetry. Original
artifacts remain the provenance for the finite exact calculations;
the printed formulas and tables make their contents available for
inspection within this expanded draft.
